\documentclass[12pt,a4paper,oneside]{report}

% --- Motor: LuaLaTeX. ELIMINADOS inputenc y fontenc (eran de pdfLaTeX) ---
\usepackage{fontspec}                 % requiere LuaLaTeX/XeLaTeX
\defaultfontfeatures{Ligatures=TeX}
\setmainfont{TeX Gyre Termes}         % clon libre de Times (estilo LNCS)
\setsansfont{TeX Gyre Heros}          % clon de Helvetica (titulares/sans)
\setmonofont{Latin Modern Mono}[Scale=MatchLowercase]  % código

\usepackage[spanish,es-noindentfirst,es-noshorthands]{babel}

% --- Márgenes "página completa" (mancha ~160 mm). Explícito en geometry ---
\usepackage[a4paper,top=2.5cm,bottom=2.5cm,left=2.5cm,right=2.5cm]{geometry}
% Variante más "llena" (mancha 170 mm) si la guía lo exige: left=right=2cm.

% --- Gráficos, tablas, color ---------------------------------------
\usepackage{graphicx}
\graphicspath{{figuras/img/}}
\usepackage{booktabs}
\usepackage{array}
\usepackage{longtable}
\usepackage{tabularx}
\usepackage{xcolor}
\usepackage{float}

\usepackage{pgfplots}                 % (carga tikz automáticamente)
\pgfplotsset{compat=1.18}
\usepgfplotslibrary{statistics}
\usetikzlibrary{positioning, arrows.meta, shapes.geometric, shapes.misc}

% --- Código fuente Java (endurecido contra desbordes) --------------
\usepackage{listings}
\lstset{
  language=Java,
  basicstyle=\ttfamily\footnotesize,
  keywordstyle=\color{blue!60!black}\bfseries,
  commentstyle=\color{green!40!black}\itshape,
  stringstyle=\color{red!60!black},
  numbers=left,
  numberstyle=\tiny\color{gray},
  frame=single,
  breaklines=true,
  breakatwhitespace=false,
  breakindent=0pt,
  postbreak=\mbox{\textcolor{gray}{$\hookrightarrow$}\space},
  showstringspaces=false,
  tabsize=2,
  captionpos=b
}

% --- Nombres de captions (es): Tabla en vez de Cuadro; Listado para código ---
\addto\captionsspanish{\renewcommand{\tablename}{Tabla}}
\renewcommand{\lstlistingname}{Listado}

% --- Matemáticas y enlaces -----------------------------------------
\usepackage{amsmath}
% Math en Times (OPCIONAL; debe ir DESPUÉS de amsmath; ver advertencias):
% \usepackage{unicode-math}
% \setmathfont{TeX Gyre Termes Math}
\usepackage{microtype}
\setlength{\emergencystretch}{3em}
\usepackage{pdfpages}
\usepackage{hyperref}
\usepackage{xurl}   % permite cortar URLs/DOIs en cualquier punto, no solo en guiones
\hypersetup{hidelinks}

% --- Índice y numeración -------------------------------------------
\setcounter{tocdepth}{2}
\setcounter{secnumdepth}{3}

% --- Títulos al gusto LNCS (Times negrita, número sin "Capítulo") --
\usepackage{titlesec}
\titleformat{\chapter}[hang]
  {\normalfont\huge\bfseries}{\thechapter}{1ex}{}
\titlespacing*{\chapter}{0pt}{0pt}{2ex}
\titleformat{\section}
  {\normalfont\large\bfseries}{\thesection}{1ex}{}
\titleformat{\subsection}
  {\normalfont\normalsize\bfseries}{\thesubsection}{1ex}{}

% =====================================================================
%  Marcadores reservados de trazabilidad (NO eliminar)
%  Hacen visibles y localizables los pendientes exigidos por las
%  instrucciones del repositorio (AGENTS.md, CLAUDE.md).
% =====================================================================
\newcommand{\pendientecita}[1][]{{\color{red}\textbf{[PENDIENTE DE CITA#1]}}}
\newcommand{\pendienteverif}[1][]{{\color{orange!80!black}\textbf{[PENDIENTE DE VERIFICACIÓN#1]}}}
\newcommand{\pendientemetodo}[1][]{{\color{violet}\textbf{[PENDIENTE DE REDACCIÓN METODOLÓGICA#1]}}}
\newcommand{\pendienteredaccion}[1][]{{\color{gray}\textbf{[PENDIENTE DE REDACCIÓN#1]}}}
\newcommand{\hipotesis}[1][]{{\color{blue!70!black}\textbf{[HIPÓTESIS DE TRABAJO#1]}}}

% =====================================================================
\begin{document}

% --- Portada --------------------------------------------------------
% Sustituye portada.pdf por el nombre real del fichero una vez lo tengas.
% Colócalo en la misma carpeta que main.tex (memoria/).
% Si el PDF de portada ocupa más de una página, cambia pages=1 por pages={1,2}.
\includepdf[pages=1, fitpaper=true]{Portada_TFM.pdf}

% --- Declaración sobre el uso de IA (revisar y firmar) -------------
% =====================================================================
%  Declaración sobre el uso de inteligencia artificial.
%  BORRADOR: el autor debe revisarlo, ajustarlo a su uso REAL y firmarlo.
%  Sustituir [Nombre y apellidos del autor] por el nombre real.
%  Front matter sin numerar.
% =====================================================================

\chapter*{Declaración sobre el uso de inteligencia artificial}
\addcontentsline{toc}{chapter}{Declaración sobre el uso de inteligencia artificial}

En cumplimiento de la normativa de la Universidad de Málaga sobre el uso de
herramientas de inteligencia artificial (IA) en los trabajos académicos, el
autor de esta memoria declara de forma expresa lo siguiente:

\begin{itemize}
  \item Durante la realización de este Trabajo Fin de Máster se han empleado
  herramientas de IA generativa basadas en grandes modelos de lenguaje como
  \emph{apoyo}, y su uso se declara abiertamente en este apartado.

  \item Dichas herramientas se han utilizado, en concreto, para: la asistencia
  en la redacción y la revisión de estilo del texto en \LaTeX; el apoyo en la
  implementación, la depuración y la refactorización del código del prototipo; y
  la ayuda en la organización y la revisión del documento.

  \item El uso se ha limitado a tareas de asistencia. Las preguntas de
  investigación, el diseño del enfoque, las decisiones metodológicas, los
  experimentos y sus resultados, y las conclusiones son obra del autor.

  \item Ningún resultado experimental, dato ni referencia bibliográfica ha sido
  generado o inventado por la IA: todos los resultados proceden de la ejecución
  reproducible del prototipo (\autoref{ap:replicacion}) y todas las referencias
  han sido verificadas por el autor.

  \item El autor ha revisado, comprendido y validado la totalidad del contenido
  presentado, y asume la plena responsabilidad sobre él.
\end{itemize}

\vspace{2em}
\noindent Málaga, \today

\vspace{3em}
\noindent Fdo.: [Nombre y apellidos del autor]


% --- Resumen / Abstract / Keywords ---------------------------------
% =====================================================================
%  Resumen / Abstract / Keywords  (Bloque UMA obligatorio)
%  ESTADO: REDACTADO. Cifras verificadas contra output/ (2026-06-03).
%  Front matter sin numerar (report).
% =====================================================================

\chapter*{Resumen}
\addcontentsline{toc}{chapter}{Resumen}

El mantenimiento del software destina una fracción muy elevada del tiempo de
desarrollo a la comprensión del código. La complejidad cognitiva, métrica
propuesta por SonarSource, cuantifica ese esfuerzo penalizando especialmente el
anidamiento de estructuras de control. Los condicionales anidados son uno de sus
principales contribuyentes y, en muchos casos, pueden combinarse de forma segura
en una única condición conjuntiva, reduciendo así la complejidad sin alterar el
comportamiento observable.

Este trabajo formaliza las condiciones bajo las cuales esa combinación es segura
(un conjunto de cuatro precondiciones estructurales verificables sobre el árbol
de sintaxis abstracta) e implementa un prototipo determinista que las comprueba y aplica la
transformación. El método experimental se articula en torno a tres preguntas de
investigación: qué casos son combinables (RQ1), qué impacto tiene la
transformación en la complejidad cognitiva (RQ2) y si los grandes modelos de
lenguaje pueden realizarla de forma correcta y automática (RQ3).

Sobre un corpus de 126 métodos Java de proyectos de código abierto, el prototipo
identifica 111 casos elegibles y reduce la complejidad cognitiva en un total de
340 puntos (todos los casos elegibles mejoran), con una coincidencia exacta
del 100\,\% frente a SonarQube en el corpus completo. En cuanto a los modelos de
lenguaje, exhiben perfiles opuestos: uno conservador, que nunca transforma en
falso, y otro más capaz pero con algún falso positivo; ninguno iguala la
garantía ni la eficiencia del prototipo determinista. El trabajo aporta, además,
un paquete de replicación completo.

\vspace{1em}
\noindent\textbf{Palabras clave:} complejidad cognitiva, refactorización,
condicionales anidados, Java, árbol de sintaxis abstracta, JavaParser, grandes
modelos de lenguaje, Design Science Research.

\chapter*{Abstract}
\addcontentsline{toc}{chapter}{Abstract}

Software maintenance devotes a large fraction of development time to code
comprehension. Cognitive complexity, a metric proposed by SonarSource, quantifies
that effort by penalising nesting in control structures. Nested conditionals are
among its main contributors, and in many cases they can be safely combined into a
single conjunctive condition, reducing complexity without altering observable
behaviour.

This work formalises the conditions under which such combination is safe (a set
of four structural verifiable preconditions on the abstract syntax tree) and implements a
deterministic prototype that checks them and applies the transformation. The
experimental method addresses three research questions: which cases are
combinable (RQ1), what impact the transformation has on cognitive complexity
(RQ2), and whether large language models can perform it correctly and
automatically (RQ3).

On a corpus of 126 Java methods from open-source projects, the prototype
identifies 111 eligible cases and reduces cognitive complexity by a total of 340
points (all eligible cases improve), with 100\,\% exact agreement against
SonarQube on the full corpus. Regarding the language models, they exhibit
opposite profiles: one conservative, never producing a false transformation, and
one more capable but with occasional false positives; neither matches the
correctness guarantees or the efficiency of the deterministic prototype. The work
also delivers a complete replication package.

\vspace{1em}
\noindent\textbf{Keywords:} cognitive complexity, refactoring, nested
conditionals, Java, abstract syntax tree, JavaParser, large language models,
Design Science Research.


% --- Índice ---------------------------------------------------------
% Con `report` el índice es nativo. Recordatorio: se rellena en la 2ª
% pasada de compilación (usar latexmk -pdf para evitar el problema).
\tableofcontents
\clearpage

% --- Cuerpo de la memoria ------------------------------------------
% =====================================================================
%  Capítulo 1. Introducción  (Bloque UMA obligatorio)
%  ESTADO: andamiaje + objetivos/contribuciones estables. La redacción
%  definitiva de §1.1 y §1.6 se cerrará tras consolidar resultados.
% =====================================================================
\chapter{Introducción}
\label{cap:introduccion}

% --- Objetivo del capítulo ------------------------------------------
% Presentar contexto, motivación, objetivos, preguntas de investigación
% y estructura del documento.

\section{Contexto y motivación}
\label{sec:contexto}
\pendienteredaccion{} La complejidad cognitiva es una métrica de
mantenibilidad que estima el esfuerzo humano necesario para comprender un
fragmento de código \pendientecita{}. A diferencia de la complejidad
ciclomática, penaliza especialmente el anidamiento de estructuras de
control, lo que convierte a los condicionales anidados en uno de sus
principales contribuyentes \pendientecita{}. Esta sección situará la
frecuencia de dichos condicionales en código real (apoyándose en el
escaneo masivo descrito en la \autoref{sec:corpus}) y motivará por qué
su simplificación es relevante para la calidad del software.

\section{Problema abordado}
\label{sec:problema-intro}
El trabajo aborda la \emph{combinación segura de sentencias condicionales
anidadas} en Java, con foco en el patrón canónico
\lstinline|if (A) { if (B) { S } }| $\rightarrow$ \lstinline|if (A && B) { S }|.
La descripción detallada del problema se ofrece en el
\autoref{cap:problema}.

\section{Objetivos}
\label{sec:objetivos-intro}
\paragraph{Objetivo general.}
Diseñar, implementar y evaluar un enfoque para detectar y refactorizar
sentencias condicionales anidadas en Java con el fin de reducir la
complejidad cognitiva del código, preservando su semántica cuando
corresponda; y comparar este enfoque con refactorizaciones generadas por
grandes modelos de lenguaje bajo un protocolo experimental controlado.

\paragraph{Objetivos específicos.}
\begin{enumerate}
  \item Identificar y formalizar los casos en que es posible combinar
  condicionales anidados de forma segura, documentando precondiciones de
  aplicabilidad y resultado de la transformación.
  \item Diseñar e implementar un prototipo determinista que detecte dichas
  oportunidades sobre el árbol de sintaxis abstracta (AST) y aplique la
  transformación.
  \item Medir el impacto de las refactorizaciones sobre la complejidad
  cognitiva con una métrica definida y reproducible.
  \item Comparar las refactorizaciones del prototipo con las producidas por
  grandes modelos de lenguaje, bajo protocolo controlado, en términos de
  corrección, completitud y calidad.
\end{enumerate}

\section{Preguntas de investigación}
\label{sec:rq-intro}
% IMPORTANTE: RQ oficiales. Reproducir SIEMPRE literalmente. No reformular.
\begin{description}
  \item[RQ1.] ¿En qué casos se pueden combinar sentencias condicionales anidadas y cuál sería el resultado?
  \item[RQ2.] ¿Qué impacto tienen estas refactorizaciones en la complejidad cognitiva?
  \item[RQ3.] ¿Pueden los grandes modelos de lenguaje realizar esta refactorización de forma correcta y automática?
\end{description}
Su operacionalización (sin reformular el texto oficial) se detalla en la
\autoref{sec:operacionalizacion}.

\section{Contribuciones del TFM}
\label{sec:contribuciones}
\begin{enumerate}
  \item La formalización de un conjunto cerrado de precondiciones
  (P1--P5) para la combinación segura de \lstinline|if| anidados, con su
  catálogo de motivos de descarte trazable.
  \item Un prototipo determinista basado en AST con un pipeline
  experimental reproducible (carga de corpus, detección, transformación,
  medición y exportación de artefactos).
  \item Una comparación controlada del prototipo frente a LLMs bajo un
  protocolo congelado y un oráculo de validación automático.
  \item Un paquete de replicación con corpus, scripts y artefactos
  versionados.
\end{enumerate}

\section{Estructura del documento}
\label{sec:estructura}
\pendienteredaccion{} El \autoref{cap:antecedentes} presenta los
antecedentes; el \autoref{cap:problema} describe el problema y
operacionaliza las RQ; el \autoref{cap:propuesta} detalla la propuesta
(diseño, arquitectura, implementación, diseño experimental, resultados y
amenazas a la validez); y el \autoref{cap:conclusiones} recoge
conclusiones y trabajo futuro. (Texto a afinar al cierre.)


  \chapter{Antecedentes}
  \label{cap:antecedentes}


  \section{Complejidad cognitiva}
  \label{sec:cc-antecedentes}

  La medición de la complejidad del software tiene un antecedente clásico en la
  \emph{complejidad ciclomática} de McCabe~\cite{10.1109/TSE.1976.233837}, que
  cuenta el número de caminos linealmente independientes del grafo de control de
  un método. Se trata de una métrica útil para estimar el esfuerzo de prueba, pero
  no fue concebida para medir lo difícil que resulta comprender el código. En particular, trata de forma parecida estructuras que un lector humano no interpreta del mismo modo
  y no penaliza de forma explícita el anidamiento, que suele
  ser uno de los factores más costosos de seguir mentalmente.

  Para cubrir esa carencia, SonarSource propuso la \emph{complejidad
  cognitiva}~\cite{10.1145/3194164.3194186}, una métrica diseñada
  específicamente para aproximar el esfuerzo de comprensión. Su modelo se apoya
  en tres tipos de incremento: un incremento estructural por cada ruptura del flujo
  lineal de lectura, un incremento por anidamiento que penaliza más las
  estructuras cuanto más profundas están y un incremento asociado a ciertas
  secuencias de operadores lógicos. Es precisamente este componente de
  anidamiento el que convierte a los condicionales anidados en uno de los
  principales contribuyentes a la complejidad cognitiva y el que motiva el patrón
  objetivo de este trabajo.

  El uso de la complejidad cognitiva como aproximación a la comprensibilidad
  cuenta además con respaldo empírico. Muñoz Barón et
  al.~\cite{10.1145/3382494.3410636} encuentran, sobre un conjunto amplio de
  datos de comprensión, una correlación significativa entre esta métrica y el
  tiempo que tardan los desarrolladores en entender el código. Esto respalda su
  uso como criterio de mejora en RQ2.

  Conviene matizar, no obstante, que la métrica empleada en este trabajo no es la
  implementación oficial de SonarQube, sino una implementación \emph{proxy} del
  modelo (\autoref{sec:proxy}). Su validez de constructo se acota mediante el
  contraste descrito en la \autoref{sec:sonar-val}.

  \section{Refactorización y preservación semántica}
  \label{sec:refactor-antecedentes}

  Se entiende por refactorización el conjunto de transformaciones que mejoran
  la estructura interna del código sin alterar su comportamiento observable.
  La referencia clásica en este ámbito es Fowler~\cite{fowler2018refactoring},
  que organiza la disciplina en torno a un catálogo de transformaciones pequeñas
  y con nombre propio, aplicadas de forma incremental. Entre ellas se encuentra
  la consolidación de expresiones condicionales, que es precisamente la familia a
  la que pertenece el patrón estudiado en este trabajo: fusionar condicionales
  anidados en una única condición conjuntiva.

  La idea de que una refactorización debe preservar el comportamiento se remonta
  a la tesis de Opdyke~\cite{opdyke1992refactoring}, donde estas transformaciones
  se definen como operaciones que conservan el comportamiento del programa
  \emph{siempre que se cumplan sus precondiciones}. Este principio es el que
  guía el enfoque del prototipo. En lugar de asumir que una transformación es
  segura en abstracto, su aplicación se restringe a aquellos casos en los que un
  conjunto de precondiciones verificables, las condiciones P1--P5 de la
  \autoref{sec:precondiciones}, permite garantizar la equivalencia.

  La noción de equivalencia que se adopta en este trabajo es la de equivalencia
  observable: el código refactorizado debe producir el mismo efecto que el
  original para cualquier entrada. En el caso del patrón
  \lstinline|if(A){if(B)}| frente a \lstinline|if(A && B)|, esta equivalencia
  depende de forma crítica de que las condiciones no introduzcan efectos
  colaterales ni alteren la semántica de cortocircuito.

  El trabajo de Mens y Tourwé~\cite{10.1109/TSE.2004.1265817} ofrece una visión
  general de la disciplina y distingue varias actividades dentro del proceso de
  refactorización: identificar dónde refactorizar, comprobar que la
  transformación preserva el comportamiento, aplicarla y evaluar su efecto sobre
  la calidad del software. Esta descomposición encaja con la estructura del
  prototipo desarrollado en este TFM, que se organiza precisamente en torno a
  detección, verificación de precondiciones, transformación y medición del
  impacto.
  La
  \autoref{tab:actividades-pipeline} hace explícita esta correspondencia.

  \begin{table}[ht]
  \centering
  \caption{Correspondencia entre las actividades de refactorización descritas
  por Mens y Tourwé~\cite{10.1109/TSE.2004.1265817} y las etapas del prototipo.}
  \label{tab:actividades-pipeline}
  \begin{tabular}{@{}lll@{}}
  \toprule
  \textbf{Actividad (Mens y Tourwé)} & \textbf{Etapa del prototipo} & \textbf{Componente} \\
  \midrule
  Identificar dónde refactorizar   & Detección del patrón      & \texttt{NestedIfDetector} \\
  Garantizar la preservación       & Verificación de P1--P5    & Precondiciones \\
  Aplicar la transformación        & Refactorización           & \texttt{NestedIfTransformer} \\
  Evaluar el efecto                & Medición del impacto      & CC \emph{proxy} \\
  \bottomrule
  \end{tabular}
  \end{table}

  \section{Análisis estático sobre AST con JavaParser}
  \label{sec:ast-antecedentes}

  Detectar y transformar el patrón objetivo exige operar sobre la estructura
  sintáctica del código, no sobre su representación textual. La forma habitual de abordar este tipo de problema es mediante un árbol de sintaxis abstracta (AST)
  , en el que cada construcción del lenguaje, por ejemplo, una sentencia
  \lstinline|if|, una expresión binaria o un bloque, se representa como un
  nodo con una posición bien definida dentro de la jerarquía del programa.
  Trabajar sobre el AST evita la fragilidad de los enfoques basados en texto y
  permite razonar de forma directa sobre relaciones estructurales como el
  anidamiento. La \autoref{fig:ast-anidado} ilustra esta idea sobre el patrón
  objetivo: el anidamiento sintáctico de los dos \lstinline|if| se traduce en una
  relación de descendencia entre nodos del árbol, de modo que el \lstinline|if|
  interno aparece como descendiente de la rama \lstinline|then| del externo.

  \begin{figure}[ht]
  \centering
  \includegraphics[width=0.45\linewidth]{figuras/img/AST.png}
  \caption{Árbol de sintaxis abstracta (AST) del patrón objetivo de dos
  sentencias \lstinline|if| anidadas. El anidamiento se representa como una
  relación de descendencia: el \lstinline|if| interno cuelga de la rama
  \lstinline|then| del externo, y el cuerpo \lstinline|S| queda en el nivel más
  profundo.}
  \label{fig:ast-anidado}
  \end{figure}

  El prototipo desarrollado en este trabajo se apoya en
  JavaParser~\cite{smith2023javaparser}, una biblioteca que permite construir el
  AST de código fuente Java y recorrerlo o modificarlo programáticamente. La
  versión concreta utilizada se fija en \texttt{pom.xml}
  (\autoref{sec:experimento}). Sobre esta base, la fase de detección recorre el
  árbol en busca de sentencias \lstinline|if| anidadas que cumplan el patrón
  objetivo, mientras que la fase de transformación reescribe el árbol
  \emph{in situ}: combina ambas condiciones en una conjunción y reubica el cuerpo
  bajo el nuevo nodo. Esta manipulación se mantiene deliberadamente
  conservadora: no reordena expresiones ni intenta simplificarlas, y preserva la
  semántica de precedencia envolviendo en paréntesis las condiciones cuando es
  necesario, por ejemplo cuando una disyunción pasa a formar parte de una
  conjunción.

  La limitación relevante para este trabajo es que el análisis realizado es
  fundamentalmente sintáctico. El AST describe la forma del programa, pero no
  resuelve por sí mismo información semántica como los tipos, los alias o el
  flujo de datos. En consecuencia, a partir del árbol no siempre puede decidirse
  si una llamada a método introduce efectos colaterales o si dos expresiones son
  equivalentes desde el punto de vista observacional. Esta frontera justifica la
  estrategia conservadora del prototipo: en lugar de intentar un análisis
  semántico completo, la aplicabilidad se restringe mediante precondiciones
  comprobables sobre la estructura y una lista explícita de construcciones
  admitidas (\autoref{sec:precondiciones}), descartando todo aquello que no pueda
  verificarse con garantías.

  \section{Grandes modelos de lenguaje aplicados a tareas de código}
  \label{sec:llm-antecedentes}

  La aplicación de grandes modelos de lenguaje (LLM) a tareas de programación
  se consolidó con Codex~\cite{chen2021codex}, un modelo entrenado sobre código
  público que demostró una capacidad notable para sintetizar programas a partir
  de descripciones en lenguaje natural y que dio lugar a asistentes de
  programación de uso masivo. Desde entonces, el interés se ha desplazado desde
  la generación de código nuevo hacia tareas sobre código existente, entre ellas
  la refactorización, que es precisamente el escenario abordado en RQ3.

  La evidencia empírica sobre refactorización con LLM es prometedora, aunque
  todavía debe interpretarse con cautela. DePalma
  et al.~\cite{10.1016/j.eswa.2024.123602} observan, sobre un conjunto de
  fragmentos de código Java, que ChatGPT propone mejoras útiles en la mayoría de
  los casos, si bien con resultados desiguales según el atributo de calidad
  considerado y sin una garantía sistemática de preservación del comportamiento.
  Liu et al.~\cite{liu2024llmrefactoring}, por su parte, encuentran que los
  modelos identifican por sí solos una fracción reducida de las oportunidades
  reales de refactorización y que su rendimiento depende en gran medida del modo
  en que la tarea se formule y acote mediante el \emph{prompt}. En la misma línea,
  Cordeiro et al.~\cite{cordeiro2024llmrefactoring} estudian la capacidad de
  refactorización de los LLMs sobre proyectos Apache reales y observan que,
  aunque los modelos producen código con una estructura mejorada, no garantizan la
  equivalencia semántica de la transformación.

  De este panorama se derivan dos implicaciones que condicionan el diseño de
  RQ3. La primera es que la corrección de una refactorización generada por un
  LLM no puede darse por supuesta, sino que debe comprobarse frente a una
  referencia fiable. En este trabajo, esa referencia la proporcionan el
  \emph{baseline} determinista del prototipo y el oráculo automático
  (\autoref{sec:oraculo}). La segunda es que la salida de estos modelos puede variar entre ejecuciones, lo que obliga a fijar un protocolo reproducible, modelo, temperatura, versión del \emph{prompt} y número de intentos por caso; para que la evaluación resulte interpretable
  (\autoref{sec:exp-rq3}).

  \section{Síntesis crítica orientada al problema}
  \label{sec:sintesis-antecedentes}

  Los antecedentes revisados convergen en el problema que aborda este trabajo,
  pero ninguno lo cubre de forma completa. La complejidad cognitiva
  (\autoref{sec:cc-antecedentes}) aporta la métrica y su justificación
  empírica; la teoría de la refactorización
  (\autoref{sec:refactor-antecedentes}) proporciona el principio de
  preservación del comportamiento condicionado al cumplimiento de
  precondiciones; el análisis sobre AST
  (\autoref{sec:ast-antecedentes}) ofrece la base tecnológica, aunque dentro
  de una frontera esencialmente sintáctica; y la literatura sobre LLM
  (\autoref{sec:llm-antecedentes}) delimita tanto el potencial de estos
  modelos como la necesidad de validar sus salidas antes de aceptarlas como
  refactorizaciones correctas.

  El trabajo más cercano al presente TFM es el de Saborido et
  al.~\cite{10.1109/ACCESS.2022.3144743}, centrado en la reducción automática
  de la complejidad cognitiva. Su propuesta modela el problema como una
  búsqueda sobre secuencias de extracción de método (\emph{Extract Method}) y
  la evalúa sobre diez proyectos de código abierto, logrando resolver una
  fracción elevada de las incidencias de complejidad cognitiva detectadas por
  SonarQube. El presente trabajo comparte con ese enfoque el objetivo general
  de reducir la complejidad cognitiva mediante refactorización automática y se
  sitúa en el mismo contexto experimental de proyectos open source
  (\autoref{sec:corpus}). Sin embargo, difiere tanto en el tipo de
  transformación como en el alcance. Mientras que Saborido et al.
  \emph{descomponen} métodos mediante extracción, este trabajo
  \emph{combina} condicionales anidados en una única condición. Se trata, por
  tanto, de operaciones distintas pero complementarias sobre una misma
  métrica. Además, aquel trabajo no aborda la capacidad de los modelos de
  lenguaje para realizar la transformación, que es precisamente el objeto de
  RQ3.

  Conviene situar también el trabajo frente a herramientas de refactorización
  automática de uso extendido. Los entornos de desarrollo integrados (IntelliJ
  IDEA, Eclipse) ofrecen un repertorio amplio de refactorizaciones, entre ellas
  la consolidación de condicionales, pero las accionan el desarrollador
  manualmente y no utilizan la complejidad cognitiva como criterio de
  decisión. OpenRewrite es un \emph{framework} de transformación AST que
  permite definir y ejecutar \emph{recetas} automáticas con garantías formales por
  receta; sin embargo, no incluye, en su catálogo estándar, una receta orientada
  específicamente a la combinación de condicionales anidados con criterio de CC y
  comprobación de precondiciones. Las reglas automáticas de SonarQube
  (\emph{quick fixes}) se limitan a sugerencias de baja complejidad y no cubren
  este patrón. Ninguna de estas herramientas aúna los tres rasgos del presente
  trabajo: transformación formalizada con precondiciones verificables, medición
  de impacto en CC y evaluación de los LLM como agente alternativo.

  La \autoref{tab:comparativa} resume este posicionamiento frente a enfoques
  previos. El hueco que cubre este trabajo se sitúa en la intersección de tres
  rasgos que, por separado, sí aparecen en la literatura: una transformación
  concreta acotada mediante precondiciones \emph{verificables}, el uso de la
  complejidad cognitiva como criterio de impacto y una evaluación explícita de
  la capacidad de los LLM para reproducir la refactorización frente a un
  \emph{baseline} determinista.

  \begin{table}[H]
  \centering
  \caption{Posicionamiento del trabajo frente a enfoques previos de refactorización y reducción de complejidad.}
  \label{tab:comparativa}
  \begin{tabular}{@{}p{3.1cm}p{2.8cm}p{2.8cm}p{2.5cm}@{}}
  \toprule
  \textbf{Enfoque} & \textbf{Transformación} & \textbf{Garantía de seguridad} & \textbf{Papel de los LLM} \\
  \midrule
  Refactors de IDE / catálogo de Fowler~\cite{fowler2018refactoring}
    & Repertorio amplio
    & Heurística de la herramienta
    & No considerado \\
  Saborido et al.~\cite{10.1109/ACCESS.2022.3144743}
    & Extracción de método (búsqueda)
    & Refactor que preserva el comportamiento
    & No considerado \\
  LLM para refactor~\cite{10.1016/j.eswa.2024.123602, liu2024llmrefactoring}
    & Abierta
    & Sin garantía formal
    & Constituyen el objeto de análisis \\
  \textbf{Este trabajo}
    & Combinación de condicionales anidados
    & Precondiciones P1--P5 verificables
    & Evaluados en RQ3 \\
  \bottomrule
  \end{tabular}
  \end{table}

\chapter{Descripción del problema}
\label{cap:problema}

Este capítulo delimita el problema que aborda el trabajo. Primero se formula
el problema y se fija el patrón objetivo de transformación; después se define
qué se entiende exactamente por condicional anidada \emph{combinable} y se
muestran casos próximos que no lo son; por último, se sitúa el patrón dentro
de una familia más amplia de problemas y se acota el alcance del prototipo.
La metodología, la operacionalización de las preguntas de investigación y las
hipótesis de trabajo se desarrollan en el \autoref{cap:metodologia}.

\section{Formulación del problema}
\label{sec:formulacion}

Los condicionales anidados son una de las construcciones que más contribuyen
a la complejidad cognitiva de un método: cada nivel de anidamiento añade un
salto mental que el lector debe sostener~\cite{10.1145/3194164.3194186}
(\autoref{sec:cc-antecedentes}). Sin embargo, en muchos casos dos sentencias
\lstinline|if| anidadas sin ramas alternativas no expresan dos decisiones
distintas, sino una sola condición conjuntiva escrita en dos niveles. Cuando
esto ocurre, aplanarlas en \lstinline|if (A && B)| simplifica el código sin
cambiar lo que hace. La misma idea se extiende a cadenas de más de dos
\lstinline|if| anidados, que se combinan aplicando la regla por pares de forma
sucesiva (\autoref{sec:alcance}).

El problema abordado nace de esa observación y tiene dos dimensiones. La
primera consiste en determinar en qué condiciones esa combinación es
semánticamente segura y automatizable. La segunda consiste en medir su efecto
sobre la complejidad cognitiva. Interesa responder a ambas no solo para una
herramienta determinista, sino también para un gran modelo de lenguaje, lo
que permite comparar dos formas muy distintas de afrontar la misma tarea.
Estas dos dimensiones se recogen en las preguntas de investigación del
trabajo, enunciadas en la \autoref{sec:rq-intro}.

\section{Patrón objetivo}
\label{sec:patron}

La transformación canónica estudiada parte de un anidamiento simple y se
transforma en una condición conjuntiva equivalente, como ilustra la
\autoref{fig:before-after}.

\begin{figure}[ht]
\centering
\begin{minipage}[c]{0.44\linewidth}
\centering
\textbf{Antes}\par\smallskip
\begin{lstlisting}[numbers=none]
if (A) {
    if (B) {
        S
    }
}
\end{lstlisting}
\end{minipage}%
\begin{minipage}[c]{0.10\linewidth}
\centering$\Longrightarrow$
\end{minipage}%
\begin{minipage}[c]{0.44\linewidth}
\centering
\textbf{Después}\par\smallskip
\begin{lstlisting}[numbers=none]
if (A && B) {
    S
}
\end{lstlisting}
\end{minipage}
\caption{Esquema \emph{before/after} del patrón objetivo: dos \texttt{if}
anidados (izquierda) se combinan en una única condición conjuntiva
\lstinline|A && B| (derecha). La transformación preserva el orden de
evaluación y es semánticamente equivalente bajo las precondiciones P1--P5.}
\label{fig:before-after}
\end{figure}

\noindent donde \lstinline|A| y \lstinline|B| son expresiones booleanas y
\lstinline|S| un bloque de sentencias arbitrario. La equivalencia se sostiene
bajo las precondiciones de aplicabilidad (\autoref{sec:precondiciones}): el
operador \lstinline|&&| evalúa primero \lstinline|A| y, solo si se cumple,
\lstinline|B|, preservando así el orden y el cortocircuito del anidamiento
original y, con ello, el comportamiento observable.

Las condiciones \lstinline|A| y \lstinline|B| pueden ser expresiones booleanas
compuestas. En ese caso, la combinación preserva la precedencia original
añadiendo paréntesis cuando es necesario (por ejemplo, si una de ellas es una
disyunción que pasa a formar parte de la conjunción).

Para hacer tangible el efecto sobre la métrica, conviene contar la complejidad
cognitiva de ambas versiones según el modelo de SonarSource
(\autoref{sec:cc-antecedentes}). En la versión anidada, el \lstinline|if|
externo aporta un incremento estructural (+1) y el \lstinline|if| interno
aporta otro estructural más uno de anidamiento por estar un nivel más adentro
(+2): en total, 3. En la versión combinada, el único \lstinline|if| aporta +1 y
la secuencia de operadores lógicos \lstinline|&&| aporta +1: en total, 2. La
combinación ahorra, por tanto, el incremento de anidamiento del segundo
\lstinline|if|. La \autoref{fig:cc-desglose} ilustra este desglose sobre un
método concreto, anotando la contribución de cada nodo a la complejidad
cognitiva del método.

\begin{figure}[ht]
\centering
\begin{minipage}{0.9\linewidth}
\textbf{Antes (anidado): CC del método $= 3$}
\begin{lstlisting}[numbers=none]
boolean puedeEnviar(boolean pagado, boolean enStock) {
    if (pagado) {            // +1  (incremento estructural)
        if (enStock) {       // +2  (estructural +1, anidamiento +1)
            return true;
        }
    }
    return false;
}
\end{lstlisting}
\smallskip
\textbf{Después (combinado): CC del método $= 2$}
\begin{lstlisting}[numbers=none]
boolean puedeEnviar(boolean pagado, boolean enStock) {
    if (pagado && enStock) { // +1 (if) y +1 (secuencia &&)
        return true;
    }
    return false;
}
\end{lstlisting}
\end{minipage}
\caption{Desglose de la complejidad cognitiva por nodo según el modelo de
SonarSource. En la versión anidada, el \texttt{if} externo aporta $+1$ y el
\texttt{if} interno $+2$ (incremento estructural más uno de anidamiento por
estar un nivel más adentro), de modo que el método suma CC${}=3$; al combinar
las condiciones, el único \texttt{if} aporta $+1$ y la conjunción
\lstinline|&&| otro $+1$, con CC${}=2$. La combinación elimina el incremento de
anidamiento del segundo \texttt{if}.}
\label{fig:cc-desglose}
\end{figure}

\section{Qué es una condicional anidada combinable}
\label{sec:definicion-combinable}

Se denomina \emph{condicional anidada combinable} a un par de sentencias
\lstinline|if| en el que el \lstinline|if| externo no tiene rama alternativa,
su rama \lstinline|then| contiene exactamente otra sentencia \lstinline|if|, y
esta segunda tampoco tiene rama alternativa. Para un par así, sustituir el
anidamiento por \lstinline|if (A && B)| preserva el comportamiento observable,
porque \lstinline|&&| evalúa \lstinline|A| y, solo si es cierto,
\lstinline|B|: el mismo orden y el mismo cortocircuito que en la versión
anidada. Esta es la condición de combinabilidad \emph{estructural}.

El prototipo añade una restricción adicional: exige que las condiciones estén
libres de efectos colaterales. No se debe a que la combinación deje de
preservar el comportamiento en su presencia, sino a que solo bajo esa restricción puede garantizarse la
seguridad de forma automática y verificable sobre la estructura del código,
sin recurrir a un análisis semántico que excede el alcance del prototipo. El
conjunto completo de precondiciones se formaliza en la
\autoref{sec:precondiciones}.

\section{Casos que no son combinables}
\label{sec:no-combinables}

La dificultad del problema no está en el caso favorable, sino en distinguirlo
de otros muy parecidos en los que aplanar el anidamiento \emph{cambia el
comportamiento}. Dos ejemplos permiten verlo con claridad.

Si el \lstinline|if| externo tiene una rama \lstinline|else|, esta cubre el
caso en que \lstinline|A| es falso; combinar las condiciones la haría
aplicable también cuando \lstinline|A| es cierto pero \lstinline|B| es falso,
alterando el flujo:

\begin{lstlisting}[caption={No combinable: \texttt{else} en el \texttt{if} externo.},label={lst:no-else}]
if (A) {
    if (B) { S }
} else {
    T          // se ejecuta cuando !A; "if (A && B) ... else T" lo rompe
}
\end{lstlisting}

De forma análoga, si el bloque externo contiene alguna sentencia además del
\lstinline|if| interno, esa sentencia se ejecuta siempre que se cumpla
\lstinline|A|; tras la combinación pasaría a ejecutarse solo cuando además se
cumpla \lstinline|B|:

\begin{lstlisting}[caption={No combinable: sentencia adicional en el bloque externo.},label={lst:no-extra}]
if (A) {
    T;             // se ejecuta cuando A, con independencia de B
    if (B) { S }
}
\end{lstlisting}

\noindent Estos casos no son anecdóticos: son precisamente los que obligan a
definir precondiciones explícitas y verificables en lugar de aplicar la
transformación de forma indiscriminada.

\section{Familia de problemas y alcance}
\label{sec:alcance}

La unidad de transformación es un par de \lstinline|if| anidados. Una cadena de
$N$ \lstinline|if| puramente anidados (cada nivel un único \lstinline|if|, sin
\lstinline|else| ni sentencias intercaladas) se reduce aplicando la regla de
combinación de forma repetida, un par cada vez, hasta que una precondición lo
impida; por eso una cadena puede combinarse total o parcialmente: si un nivel
intermedio incumple P1--P5, la combinación se detiene en ese punto. Quedan
fuera del alcance, como extensiones naturales (\autoref{cap:conclusiones}), las
variantes que no son anidamiento puro: cadenas con \lstinline|else if|,
secuencias de \lstinline|if| hermanos (no anidados) o combinaciones que exigen
reordenar el flujo.

La \autoref{tab:alcance} resume qué entra y qué queda fuera del alcance
inicial. Las construcciones excluidas no son irrelevantes; simplemente son
aquellas cuya seguridad no puede establecerse con garantías dentro del alcance
del prototipo.

\begin{table}[ht]
\centering
\caption{Delimitación del alcance del patrón objetivo.}
\label{tab:alcance}
\begin{tabular}{@{}p{6.1cm}p{6.1cm}@{}}
\toprule
\textbf{Dentro del alcance} & \textbf{Fuera del alcance inicial} \\
\midrule
Un \lstinline|if| que contiene exactamente otro \lstinline|if| (anidamiento simple)
  & Ramas \lstinline|else| o \lstinline|else if| en cualquiera de los dos \lstinline|if| \\
Ausencia de sentencias adicionales en el bloque externo
  & Sentencias antes o después del \lstinline|if| interno \\
Condiciones \lstinline|A| y \lstinline|B| sin efectos colaterales
  & Efectos colaterales o asignaciones en las condiciones \\
Cuerpo \lstinline|S| arbitrario
  & Expresiones lambda, \emph{streams}, concurrencia \\
Combinación mediante \lstinline|&&| con preservación del cortocircuito
  & Construcciones cuya equivalencia no se garantiza sintácticamente \\
\bottomrule
\end{tabular}
\end{table}

\chapter{Metodología}
\label{cap:metodologia}

Este capítulo describe el enfoque metodológico que vertebra el trabajo. El
problema abordado consiste en proponer una transformación de código, construir
un prototipo que la realice y comprobar empíricamente sus efectos; esa
naturaleza encaja de forma natural con el paradigma de \emph{Design Science
Research} (DSR), orientado a generar conocimiento mediante la construcción y
evaluación de artefactos.

\section{Design Science Research}
\label{sec:dsr-metodo}

DSR aborda la investigación a través de un artefacto que resuelve un problema
relevante y cuya utilidad se demuestra mediante evaluación rigurosa. Hevner et
al.~\cite{10.2307/25148625} sitúan en el núcleo de este paradigma la
articulación entre construir el artefacto y evaluarlo frente al problema que
pretende resolver. Peffers et al.~\cite{10.2753/MIS0742-1222240302} concretan
ese principio en una secuencia de seis actividades: identificar el problema y
motivarlo, fijar los objetivos de la solución, diseñar y desarrollar el
artefacto, demostrarlo, evaluarlo y comunicarlo. Esta secuencia proporciona una
estructura reconocible para presentar este tipo de investigación. La
\autoref{fig:dsr-ciclo} representa ese ciclo y la realimentación entre
construcción y evaluación que lo caracteriza.

% TODO: REGENERAR figuras/img/DSR.png desde el fuente actualizado
% figuras/src/fig-dsr-ciclo.mmd (el texto "impacto en CC" se cambió a
% "impacto en complejidad cognitiva"; el PNG actual aún muestra "CC").
\begin{figure}[ht]
\centering
\includegraphics[width=0.4\linewidth]{figuras/img/DSR.png}
\caption{Ciclo de \emph{Design Science Research} seguido en el trabajo, según
las actividades de Peffers et al.~\cite{10.2753/MIS0742-1222240302}. La
construcción del artefacto y su evaluación se articulan en un ciclo con
realimentación; las anotaciones indican su materialización en este TFM.}
\label{fig:dsr-ciclo}
\end{figure}

El artefacto de este trabajo es un prototipo determinista de detección y
refactorización de condicionales anidados. Su carácter determinista cumple
además una función metodológica de segundo orden: actúa como \emph{baseline}
verificable, es decir, como la referencia fiable contra la que se contrastan
después las refactorizaciones producidas por los grandes modelos de lenguaje
en RQ3. De este modo, la doble naturaleza de construir y comprobar, propia de
DSR, no solo organiza la propuesta, sino que sostiene el diseño experimental.



\section{Operacionalización en el TFM}
\label{sec:dsr-operacion}
\label{sec:operacionalizacion}

La evaluación del artefacto se articula en torno a las tres preguntas de
investigación: qué casos de condicionales anidados son combinables y con qué
resultado (RQ1), qué efecto tiene la transformación sobre la complejidad
cognitiva (RQ2) y en qué medida un modelo de lenguaje reproduce la
transformación de forma correcta y automática (RQ3). La
\autoref{tab:dsrm-tfm} relaciona las actividades de la metodología de Peffers
et al. con su materialización concreta en este trabajo y con los capítulos que
las desarrollan.

\begin{table}[ht]
\centering
\caption{Actividades de la metodología DSR de Peffers et
al.~\cite{10.2753/MIS0742-1222240302} y su materialización en el trabajo.}
\label{tab:dsrm-tfm}
\begin{tabular}{@{}p{4.2cm}p{5.2cm}p{2.6cm}@{}}
\toprule
\textbf{Actividad (Peffers et al.)} & \textbf{Materialización en el trabajo} & \textbf{Capítulo} \\
\midrule
Identificar el problema y motivarlo & Coste cognitivo del anidamiento de condicionales & \ref{cap:problema} \\
Fijar los objetivos de la solución  & Objetivos y preguntas de investigación RQ1--RQ3 & \ref{cap:introduccion} \\
Diseñar y desarrollar el artefacto  & Precondiciones P1--P4, prototipo basado en AST y oráculo & \ref{cap:propuesta} \\
Demostrar el artefacto              & Ejecución sobre el corpus piloto y real & \ref{sec:experimento} \\
Evaluar                             & Medición de $\Delta$ de complejidad y evaluación de los LLMs & \ref{sec:resultados} \\
Comunicar                           & La presente memoria y el paquete de replicación & --- \\
\bottomrule
\end{tabular}
\end{table}

La parte de \emph{construcción} comprende la formalización de las
precondiciones de aplicabilidad y el desarrollo del prototipo
(\autoref{cap:propuesta}). La parte de \emph{evaluación} se apoya en
protocolos deliberadamente congelados (con todos sus parámetros fijados de
antemano y registrados, de modo que no varíen entre ejecuciones) y en
artefactos persistentes y auditables, de modo que el procedimiento pueda
repetirse y conducir al mismo resultado; sus detalles se describen en el diseño experimental
(\autoref{sec:experimento}).

Con ese marco, cada pregunta se operacionaliza, sin alterar su enunciado
(\autoref{sec:rq-intro}), en los términos siguientes.

\paragraph{RQ1 — \textquestiondown{}En qué casos se pueden combinar sentencias condicionales anidadas y cuál sería el resultado?}
Se aborda de forma constructiva: se delimita el patrón objetivo y se formalizan
las precondiciones que hacen segura su combinación, que el detector aplica sobre
el corpus para separar los casos elegibles de los que no lo son. En concreto:
\begin{itemize}
  \item Se estudian condicionales anidados de tipo \lstinline|if| en Java.
  \item Foco inicial: un \lstinline|if| que contiene exactamente otro
  \lstinline|if|, sin ramas \lstinline|else| ni \lstinline|else if|, y sin
  sentencias adicionales en el bloque externo.
  \item El resultado es la transformación
  \lstinline|if (A) { if (B) { S } }| $\rightarrow$ \lstinline|if (A && B) { S }|.
  \item Se documentan las precondiciones de aplicabilidad
  (\autoref{sec:precondiciones}).
\end{itemize}

\paragraph{RQ2 — \textquestiondown{}Qué impacto tienen estas refactorizaciones en la complejidad cognitiva?}
Se aborda de forma empírica: se mide la complejidad cognitiva de cada método
antes y después de aplicar la transformación y se agrega el efecto sobre todo el
corpus. En concreto:
\begin{itemize}
  \item La complejidad cognitiva se mide siguiendo el modelo de
  SonarSource~\cite{10.1145/3194164.3194186}, mediante una implementación
  \emph{proxy} cuyas limitaciones se declaran en la \autoref{sec:proxy}.
  \item Se compara el valor antes y después de la refactorización en cada caso.
  \item Se reporta la reducción absoluta y relativa, por caso y agregada.
\end{itemize}

\paragraph{RQ3 — \textquestiondown{}Pueden los grandes modelos de lenguaje realizar esta refactorización de forma correcta y automática?}
Se aborda de forma comparativa: se solicita la misma transformación a varios
modelos de lenguaje bajo un protocolo controlado y sus salidas se contrastan con
el \emph{baseline} determinista mediante un oráculo automático. En concreto:
\begin{itemize}
  \item Se evalúa sobre el corpus de evaluación (\autoref{sec:corpus}).
  \item Se solicita a uno o más LLMs la misma refactorización bajo un
  \emph{prompt} estandarizado y un protocolo congelado.
  \item Se comparan los resultados en corrección semántica, completitud de
  detección y calidad del código, mediante un oráculo automático
  (\autoref{fig:oraculo}).
\end{itemize}

\paragraph{Criterio de corrección.}
Una transformación se considera correcta si preserva el comportamiento
observable del método. Esa corrección se sostiene en dos niveles. Por
construcción: bajo las precondiciones de aplicabilidad
(\autoref{sec:precondiciones}), el operador \lstinline|&&| respeta el orden y
el cortocircuito de la evaluación original, de modo que ambas versiones
producen el mismo efecto para toda entrada. Y de forma dinámica: el resultado
de cada caso elegible se compila con el compilador de Java para comprobar que
la transformación no genera código inválido. La ejecución de las pruebas de
comportamiento de los proyectos de origen excede el alcance del prototipo,
pues requeriría reconstruir y configurar cada proyecto, y se señala como
trabajo futuro (\autoref{cap:conclusiones}).

\section{Hipótesis de trabajo}
\label{sec:hipotesis}
El trabajo se guía por tres hipótesis que se contrastan con los resultados y
que, deliberadamente, no anticipan cifras.

\begin{description}
  \item[H1.] Existe un subconjunto no trivial de condicionales anidados en
  código Java real para los que la combinación es segura bajo las
  precondiciones de aplicabilidad.
  \item[H2.] La combinación reduce la complejidad cognitiva (proxy) de forma
  medible en los casos elegibles.
  \item[H3.] Los LLMs pueden reproducir la transformación con tasas de acierto
  elevadas en casos elegibles, pero pueden fallar en el rechazo de casos no
  elegibles.
\end{description}


\chapter{Detalles de la propuesta}
\label{cap:propuesta}

Este capítulo desarrolla la propuesta del trabajo: un prototipo determinista que
detecta y combina condicionales anidados de forma segura. Se presenta primero
una visión general del enfoque y su encaje en la metodología \emph{Design
Science Research}; a continuación, el diseño del enfoque, con el patrón objetivo
y las precondiciones de aplicabilidad que delimitan cuándo la transformación es
segura; después, la arquitectura del prototipo en paquetes y flujos de datos; y,
por último, los detalles de implementación de cada componente.

\section{Visión general}
\label{sec:vision-general}
La propuesta se plantea como un proyecto de \emph{Design Science Research}
(DSR)~\cite{10.2307/25148625, 10.2753/MIS0742-1222240302}, porque su
estructura encaja de forma natural con el
objetivo: construir un artefacto y evaluarlo. El artefacto es un prototipo
determinista de detección y refactorización; la evaluación consiste en
medir cómo responde a cada una de las preguntas de investigación. Esta
doble naturaleza (construir y comprobar) recorre todo el trabajo. Conviene
subrayar, además, que el prototipo no es solo el objeto de estudio: actúa
como \emph{baseline verificable}, es decir, como la referencia fiable
contra la que se contrastan después las refactorizaciones de los grandes
modelos de lenguaje en RQ3.

\paragraph{Mapa RQ $\leftrightarrow$ secciones.}
La \autoref{tab:mapa-rq} relaciona cada pregunta de investigación con las
secciones que la abordan.

\begin{table}[ht]
\centering
\caption{Trazabilidad entre preguntas de investigación y secciones.}
\label{tab:mapa-rq}
\begin{tabular}{@{}lll@{}}
\toprule
\textbf{RQ} & \textbf{Aborda} & \textbf{Secciones} \\
\midrule
RQ1 & Casos combinables y resultado & \ref{sec:diseno}, \ref{sec:precondiciones}, \ref{sec:resultados-rq1} \\
RQ2 & Impacto en complejidad cognitiva & \ref{sec:proxy}, \ref{sec:exp-rq2}, \ref{sec:resultados-rq2} \\
RQ3 & LLMs: corrección y automatización & \ref{sec:llm-subsistema}, \ref{sec:exp-rq3}, \ref{sec:resultados-rq3} \\
\bottomrule
\end{tabular}
\end{table}




\section{Diseño del enfoque}
\label{sec:diseno}

Antes de entrar en los detalles, conviene fijar la lógica que guía el
enfoque. En esta sección se presenta la estrategia general y los
principios que la sostienen; a continuación, el patrón objetivo y las
precondiciones que delimitan cuándo es seguro aplicarlo; y, por último, la
política de transformación y la política iterativa con las que el
prototipo opera sobre el código.

\subsection{Estrategia general}
El enfoque responde a \textbf{RQ1} mediante un prototipo determinista que
trabaja directamente sobre el árbol de sintaxis abstracta (AST) de Java.
La decisión de partida es deliberadamente \textbf{conservadora}: ante la
menor duda sobre la equivalencia semántica de una transformación, el caso
se rechaza. Es preferible perder oportunidades de refactorización antes que
arriesgar un cambio de comportamiento, porque el prototipo aspira a servir
como baseline fiable frente al cual medir a los LLMs en RQ3.

Esta estrategia se concreta en tres principios:

\begin{enumerate}
  \item \textbf{Aplicabilidad explícita.} Toda transformación queda
  condicionada a un conjunto cerrado de precondiciones formales, evaluadas
  sobre el AST. Si alguna no se cumple, el caso se descarta con un motivo
  trazable, lo que permite justificar después por qué un fragmento no se
  combinó.
  \item \textbf{Preservación del orden de evaluación.} La condición
  combinada respeta el orden original \lstinline|A && B|. No se reordenan
  los operandos ni se simplifica el álgebra booleana, de modo que la
  evaluación en cortocircuito se comporta igual que antes.
  \item \textbf{Determinismo y reproducibilidad.} Para una misma entrada,
  detector y transformador producen siempre la misma salida. Es un
  requisito imprescindible para que el contraste con los LLMs en RQ3 sea
  justo y repetible.
\end{enumerate}

En conjunto, estos principios trazan una frontera nítida entre lo que el
prototipo se atreve a transformar y lo que prefiere dejar intacto, frontera
que las precondiciones de la \autoref{sec:precondiciones} formalizan.

\subsection{Patrón objetivo de transformación}
El patrón canónico (\autoref{fig:before-after}) es semánticamente equivalente
al original \emph{bajo} las precondiciones estructurales P1--P4. La equivalencia
se sostiene en el cortocircuito de \lstinline|&&|~\cite{jls17}: como \lstinline|A && B|
equivale a \lstinline|A ? B : false|, evalúa \lstinline|A| y, solo si es cierto,
\lstinline|B|, exactamente igual que el anidamiento. Por ello la combinación
preserva el comportamiento observable \emph{con independencia de que las
condiciones tengan efectos colaterales}: estos ocurren en el mismo orden y bajo
las mismas condiciones en ambas versiones. Las precondiciones que realmente
importan son, por tanto, las puramente estructurales (P1--P4); no se requiere
análisis de pureza ni de flujo de datos para garantizar la corrección.

\subsection{Precondiciones de aplicabilidad (P1--P4)}
\label{sec:precondiciones}
El detector evalúa, para cada nodo \lstinline|IfStmt| del AST, cuatro
precondiciones \emph{estructurales}, necesarias y conjuntamente suficientes
para considerar el caso elegible (la razón de la suficiencia, el cortocircuito
de \lstinline|&&|, se explica en la \autoref{sec:patron}). Los nombres en monoespaciado como
\lstinline|IfStmt| o \lstinline|BlockStmt| son
tipos de nodo del AST de JavaParser, no clases del prototipo; las clases y
enumeraciones propias del prototipo (incluidos los motivos de descarte) se
recogen en el diagrama de clases (\autoref{fig:clases-nucleo}).

\begin{description}
  \item[P1 — Ausencia de \texttt{else} externo.] El \lstinline|if| externo
  no tiene rama \lstinline|else| ni \lstinline|else if|. Un flujo
  alternativo no se preserva con la combinación \lstinline|A && B|.
  \item[P2 — Rama externa unitaria.] La rama \lstinline|then| del
  \lstinline|if| externo consiste en \emph{exactamente} una sentencia, ya
  sea encerrada en un bloque \lstinline|{...}| de una sola sentencia o
  escrita directamente sin llaves.
  \item[P3 — La sentencia única es un \texttt{if}.] La única sentencia de la
  rama \lstinline|then| exigida por P2 debe ser, a su vez, otro
  \lstinline|IfStmt| (el \lstinline|if| interno candidato). El detector
  contempla ambas formas, con y sin llaves (p.\,ej.\
  \lstinline|if (A) if (B) S;|).
  \item[P4 — Ausencia de \texttt{else} interno.] El \lstinline|if| interno
  tampoco tiene rama \lstinline|else| ni \lstinline|else if|.
\end{description}

En conjunto, P1 y P4 exigen que ninguno de los dos \lstinline|if| (ni el
externo ni el interno) tenga rama \lstinline|else| o \lstinline|else if| ---que
en el AST de JavaParser no es una construcción propia, sino un \lstinline|else|
cuyo cuerpo es a su vez un \lstinline|IfStmt|, de modo que la comprobación de
ausencia de \lstinline|else| (\lstinline|hasElseBranch()|) cubre ambos casos por
igual---, y P2 y P3, que el bloque externo contenga exactamente el
\lstinline|if| interno y nada más. Estas cuatro condiciones son puramente estructurales y, por el
cortocircuito de \lstinline|&&| (\autoref{sec:patron}), necesarias y
suficientes para que la combinación preserve el comportamiento.

La \autoref{fig:arbol-precondiciones} resume el flujo de decisión completo: las
precondiciones se evalúan en orden y, en cuanto una falla, el candidato se
descarta con un motivo trazable; solo si todas se cumplen el caso se declara
elegible.

\begin{figure}[ht]
\centering
\begin{tikzpicture}[
  check/.style={draw, rounded corners, fill=blue!5, align=center,
                text width=4.2cm, inner sep=4pt, font=\small},
  reject/.style={draw, fill=red!5, align=left,
                 text width=5.0cm, inner sep=3pt, font=\scriptsize},
  accept/.style={draw, thick, fill=green!12, align=center,
                 text width=4.2cm, inner sep=4pt, font=\small},
  arr/.style={-{Latex}},
  lbl/.style={font=\scriptsize, inner sep=1pt},
]
% Columna izquierda (decisiones), coordenadas explícitas
\node[check]  (p1) at (0,0)    {P1 — \textquestiondown{}\texttt{if} externo sin \texttt{else}?};
\node[check]  (p2) at (0,-2)   {P2 — \textquestiondown{}rama \texttt{then} con una sola sentencia?};
\node[check]  (p3) at (0,-4)   {P3 — \textquestiondown{}esa sentencia es un \texttt{if}?};
\node[check]  (p4) at (0,-6)   {P4 — \textquestiondown{}\texttt{if} interno sin \texttt{else}?};
\node[accept] (ok) at (0,-8)   {\textbf{Elegible:} combinar \texttt{A \&\& B}};
% Columna derecha (motivos de descarte)
\node[reject] (r1) at (8,0)    {\texttt{OUTER\_HAS\_ELSE}};
\node[reject] (r2) at (8,-2)   {\texttt{OUTER\_BLOCK\_MULTIPLE\_STATEMENTS}};
\node[reject] (r3) at (8,-4)   {\texttt{SINGLE\_STATEMENT\_NOT\_IF}};
\node[reject] (r4) at (8,-6)   {\texttt{INNER\_HAS\_ELSE}};
% Flujo "sí" (vertical)
\draw[arr] (p1) -- node[lbl,left]{sí} (p2);
\draw[arr] (p2) -- node[lbl,left]{sí} (p3);
\draw[arr] (p3) -- node[lbl,left]{sí} (p4);
\draw[arr] (p4) -- node[lbl,left]{sí} (ok);
% Flujo "no" (horizontal hacia el motivo de descarte)
\draw[arr] (p1) -- node[lbl,above]{no} (r1);
\draw[arr] (p2) -- node[lbl,above]{no} (r2);
\draw[arr] (p3) -- node[lbl,above]{no} (r3);
\draw[arr] (p4) -- node[lbl,above]{no} (r4);
\end{tikzpicture}
\caption{Árbol de decisión del detector. Las precondiciones estructurales
P1--P4 son necesarias y suficientes para la combinación; cada rama "no" produce
un motivo de descarte estable (\texttt{DiscardReason}).}
\label{fig:arbol-precondiciones}
\end{figure}

Estas precondiciones están implementadas en \texttt{NestedIfDetector} y
materializadas como categorías estables del enumerado
\texttt{DiscardReason}, que el detector emite al rechazar un candidato. El
catálogo completo de motivos (incluyendo la categoría auxiliar
\texttt{NO\_NESTED\_IF\_PATTERN}) se reproduce en el
Apéndice~D. El detector expone dos métodos de salida:
\texttt{detect(method)} (solo oportunidades aceptadas) y
\texttt{detectWithReasons(method)} (también los candidatos rechazados con
sus motivos), siendo este último el que alimenta la columna
\texttt{discardCategories} del CSV de resultados y da respuesta empírica a
RQ1.

El detector aplica estas cuatro precondiciones en el modo \texttt{STRUCTURAL}
del enumerado \texttt{DetectionMode}, que es el que sustenta los
\textbf{resultados principales} de esta memoria.

\subsection{Política de transformación y seguridad semántica}

Una vez detectada una oportunidad, \texttt{NestedIfTransformer} aplica la transformación sobre el AST siguiendo esta secuencia:

\begin{enumerate}
    \item \textbf{Guarda previa.}
    Al entrar, el transformador comprueba que ninguno de los dos \lstinline|if| tenga rama \lstinline|else|; si la tuviera, aborta devolviendo \lstinline|false| sin modificar el AST. Esta comprobación es redundante respecto a P1 y P4, pero protege frente a usos directos del transformador al margen del detector.

    \item \textbf{Clonación de las condiciones.}
    Se clonan \lstinline|A| y \lstinline|B| antes de reutilizarlas en la
    condición combinada. El \emph{aliasing} es la situación en que dos partes del
    árbol comparten el mismo nodo: como en JavaParser cada nodo tiene un único
    nodo padre, insertar directamente la condición original en el nuevo
    \lstinline|if| la desligaría de su ubicación previa y dejaría el AST en un
    estado inconsistente. Clonar produce copias independientes y evita ese efecto.

    \item \textbf{Parentizado de subexpresiones con menor precedencia.}
    Cuando alguna condición es una expresión \lstinline=||=, se encapsula mediante \lstinline|EnclosedExpr| para preservar correctamente la precedencia frente a \lstinline|&&|.

    \item \textbf{Construcción de la condición combinada.}
    Se crea un nuevo \lstinline|BinaryExpr| con operador \lstinline|AND|, manteniendo el orden original de evaluación como \lstinline|outer && inner|.

    \item \textbf{Sustitución del nodo externo.}
    La condición del \lstinline|if| externo se reemplaza por la nueva condición combinada y el cuerpo pasa a ser el del \lstinline|if| interno, previamente clonado y normalizado a \lstinline|BlockStmt|. Por último se invoca \lstinline|removeElseStmt()| como salvaguarda adicional (\emph{no-op} si, como garantiza la guarda previa, ya no hay \lstinline|else|).
\end{enumerate}
\subsection{Política iterativa una por una}
\label{sec:one-at-a-time}

Las múltiples oportunidades dentro de un mismo método se procesan \textbf{de una en una}: tras cada transformación se vuelve a ejecutar la detección sobre el AST modificado, hasta que no queden más oportunidades o se alcance un límite de seguridad de 10 pasadas (\lstinline|MAX_PASSES| en \texttt{BatchRunner}), un tope que garantiza la terminación del bucle de re-detección aunque, por un error inesperado, una pasada no llegara a reducir el número de oportunidades. Esta política se adoptó al observar que aplicar todas las oportunidades en un único pase podía producir sobreconteo cuando los nodos afectados estaban anidados unos dentro de otros.

En concreto, esta decisión garantiza que:

\begin{enumerate}
    \item \lstinline|opportunitiesApplied| refleje únicamente transformaciones realmente aplicadas;
    \item la métrica $\Delta (\Delta = after - before)$ calculada antes y después de la refactorización sea exacta para cada caso; y
    \item la trazabilidad del proceso iterativo quede registrada mediante \lstinline|totalPasses|.
\end{enumerate}

En conjunto, el diseño sacrifica cobertura a cambio de garantías: los casos con
\lstinline|else| o con sentencias adicionales en el bloque externo quedan fuera
de alcance, pues no son combinables sin alterar el comportamiento. Esta
restricción, puramente estructural, permite que el pipeline determinista actúe
como \textbf{baseline verificable} frente a los LLMs en RQ3.


\section{Arquitectura del prototipo}
\label{sec:arquitectura}

Esta sección describe la organización del prototipo: primero la división en
paquetes y responsabilidades, después el flujo de datos del pipeline
determinista (RQ1 y RQ2) y el del subsistema de evaluación de los modelos de
lenguaje (RQ3), y por último las decisiones estructurales que lo sostienen.

\subsection{Visión general}
La organización del prototipo en paquetes responde a una idea sencilla:
cada paquete asume una responsabilidad y solo una. Se parte de las tres
responsabilidades clásicas de una herramienta de análisis y
refactorización, detección, transformación y medición; y se amplían
con dos piezas propias de este trabajo: un pipeline experimental que
orquesta los \emph{lotes} (cada lote es una ejecución del proceso sobre todo
el corpus) y un subsistema dedicado a la evaluación de los LLMs
(RQ3). Esta separación no es estética; es lo que permite cambiar una
parte sin arrastrar a las demás y, sobre todo, razonar sobre cada etapa de
forma aislada. Todas las operaciones sobre código se apoyan en el AST que
proporciona la biblioteca \textbf{JavaParser}, declarada como dependencia
en \texttt{pom.xml}. La \autoref{fig:paquetes} (diagrama de paquetes) y la
\autoref{fig:pipeline-det} (flujo de datos del pipeline determinista)
acompañan esta sección.

\subsection{Paquetes y responsabilidades}
El prototipo se organiza en paquetes con responsabilidades separadas. La
\autoref{tab:paquetes-arq} resume cada paquete y sus clases principales.

\begin{longtable}{@{}p{0.30\linewidth}p{0.62\linewidth}@{}}
\caption{Paquetes del prototipo y sus clases principales.}\label{tab:paquetes-arq}\\
\toprule
\textbf{Paquete} & \textbf{Responsabilidad y clases principales} \\
\midrule
\endhead
\texttt{es.tfm.refactoring} & Punto de entrada de línea de órdenes
(\texttt{Main}). \\
\texttt{\dots detection} & Análisis del AST y evaluación de P1--P4 (y del filtro opcional P5):
\texttt{NestedIfDetector}, \texttt{RefactoringOpportunity},
\texttt{DetectionResult}, \texttt{DiscardReason}, \texttt{DetectionMode},
\texttt{MethodCallAllowlist}. \\
\texttt{\dots transformation} & Aplicación de la transformación:
\texttt{NestedIfTransformer}. \\
\texttt{\dots analysis} & Cálculo proxy de complejidad cognitiva e impacto:
\texttt{CognitiveComplexityCalculator}, \texttt{RefactoringImpact},
\texttt{RefactoringImpactAnalyzer}, \texttt{CcAnalyzeCli}. \\
\texttt{\dots validation} & Comprobación de precondiciones y validez de la
transformación: \texttt{RefactoringValidator}, \texttt{ValidationResult}. \\
\texttt{\dots scanner} & Escaneo masivo de proyectos de código abierto para
construir el corpus real: \texttt{ProjectCorpusScanner},
\texttt{CorpusScanRunner}, \texttt{CorpusFileGenerator}, \texttt{ScanFinding}. \\
\texttt{\dots experiment} & Carga de corpus, ejecución por lotes,
exportación y contraste con SonarQube: \texttt{BatchRunner},
\texttt{Rq2BatchExecutor}, \texttt{ExperimentCase},
\texttt{ExperimentResult}, \texttt{ResultExporter},
\texttt{PilotCorpusLoader}, \texttt{RealDatasetLoader},
\texttt{SonarContrastEntry}, \texttt{SonarContrastReport}. \\
\texttt{\dots llm} & Protocolo, prompt, oráculo y campaña para RQ3:
\texttt{LlmExperimentProtocol}, \texttt{LlmPromptBuilder},
\texttt{LlmEvaluationSubset}, \texttt{LlmResponseValidator},
\texttt{LlmCampaignRunner}, \texttt{LlmCampaignReport},
\texttt{LlmCampaignExporter}, \texttt{LlmResponseProvider},
\texttt{LiveLlmResponseProvider}, \texttt{PrerecordedResponseProvider},
\texttt{LlmInvocationRecord}, \texttt{LlmEvaluationResult},
\texttt{LlmVerdict}, \texttt{LlmApiConfig}, \texttt{CampaignExecutor}. \\
\bottomrule
\end{longtable}

\noindent La tabla recoge las clases principales de cada paquete; se omiten
por brevedad algunos ejecutores y variantes auxiliares (p.\,ej.\
\texttt{Rq2ComparisonExecutor}, \texttt{Rq2ScannerBatchExecutor},
\texttt{TutorCorpusLoader}, \texttt{TrapCorpusLoader} en
\texttt{experiment}, o \texttt{Rq3PromptV2Executor} y
\texttt{Rq3TrapCampaignExecutor} en \texttt{llm}), que dan soporte a
experimentos secundarios y no alteran el flujo principal descrito.

\begin{figure}[H]
\centering
\begin{tikzpicture}[
  box/.style={draw, rounded corners, fill=blue!5, align=center, font=\small},
  orch/.style={box, fill=orange!12, text width=2.6cm, inner sep=4pt, font=\small\ttfamily},
  mainbox/.style={box, fill=orange!12, text width=5.8cm, inner sep=4pt, font=\small\ttfamily},
  core/.style={box, text width=9.4cm, inner sep=6pt},
  supp/.style={box, text width=2.6cm, inner sep=4pt, font=\small\ttfamily},
  arr/.style={-{Latex}, gray!70},
]
\node[mainbox] (main) at (4,0)      {Main / ExperimentOrchestrator};
\node[orch]    (exp)  at (0.5,-2.3) {experiment};
\node[orch]    (llm)  at (4,-2.3)   {llm};
\node[orch]    (scan) at (7.5,-2.3) {scanner};
\node[core]    (core) at (4,-4.7)
  {Núcleo de análisis y transformación\\[2pt]
   \ttfamily detection \quad transformation \quad analysis};
\node[supp]    (val)  at (0.5,-6.7) {validation};
% Entrada -> orquestación
\draw[arr] (main) -- (exp);
\draw[arr] (main) -- (llm);
\draw[arr] (main) -- (scan);
% Orquestación -> núcleo (verticales, sin cruces)
\draw[arr] (exp)  -- (exp  |- core.north);
\draw[arr] (llm)  -- (llm  |- core.north);
\draw[arr] (scan) -- (scan |- core.north);
% Núcleo -> validation
\draw[arr] (core.south -| val) -- (val);
\end{tikzpicture}
\caption{Diagrama de paquetes del prototipo. La capa de entrada
(\texttt{Main}) lanza los flujos de orquestación (\texttt{experiment} para
RQ1/RQ2, \texttt{llm} para RQ3, \texttt{scanner} para construir el corpus),
que se apoyan en el núcleo de análisis y transformación
(\texttt{detection}, \texttt{transformation}, \texttt{analysis}) y en
\texttt{validation}.}
\label{fig:paquetes}
\end{figure}

La \autoref{fig:paquetes} muestra la organización en paquetes; la
\autoref{fig:clases-nucleo} desciende al nivel de clases del núcleo
determinista, con sus métodos públicos y las enumeraciones de configuración
(\texttt{DetectionMode}) y de motivos de descarte (\texttt{DiscardReason}).

% TODO: FALTA la imagen figuras/img/fig-4.4.png (solo existe el fuente .mmd).
% Mientras no se genere, el PDF muestra la ruta como texto en lugar del diagrama.
% Generarla antes de compilar la versión final con:
%   mmdc -i figuras/src/fig-4.4-clases-nucleo.mmd -o figuras/img/fig-4.4.png -w 2600
\begin{figure}[H]
\centering
\includegraphics[width=\linewidth,height=0.85\textheight,keepaspectratio]{figuras/img/fig-4.4.png}
\caption{Diagrama de clases (UML) del núcleo determinista: detección
(\texttt{NestedIfDetector}), transformación (\texttt{NestedIfTransformer}),
medición (\texttt{CognitiveComplexityCalculator}), validación
(\texttt{RefactoringValidator}) y orquestación por lotes
(\texttt{BatchRunner}), con las clases de dominio \texttt{RefactoringOpportunity}
y \texttt{DetectionResult} y las enumeraciones \texttt{DetectionMode} (modos
\texttt{STRICT}/\texttt{RELAXED}) y \texttt{DiscardReason}. Las relaciones
siguen la notación UML: agregación (rombo) del orquestador sobre los servicios
del núcleo, asociación con el modo de detección y dependencias de creación y
uso.}
\label{fig:clases-nucleo}
\end{figure}

\subsection{Flujo de datos del pipeline determinista (RQ1, RQ2)}

El flujo extremo a extremo, implementado por \texttt{BatchRunner.processCase}, se organiza en las siguientes etapas:

\begin{enumerate}
    \item \textbf{Carga del corpus.}  
    Los casos se cargan desde el \emph{classpath} y se construyen como instancias de \texttt{ExperimentCase}; cada caso es un método Java acompañado de sus metadatos (procedencia, complejidad esperada, etc.).

    \item \textbf{Parseo a AST.}  
    El código fuente se parsea mediante \lstinline|StaticJavaParser.parse(...)| y se selecciona el primer \lstinline|MethodDeclaration|. Si el parseo falla, el caso se marca como inelegible con un motivo explícito.

    \item \textbf{Detección con motivos.}  
    Se ejecuta \lstinline|detectWithReasons| para identificar oportunidades y registrar los \texttt{DiscardReason} asociados en \lstinline|discardCategories|.

    \item \textbf{Medición \emph{before}.}  
    Se calcula la complejidad cognitiva inicial del método antes de aplicar ninguna transformación.

    \item \textbf{Aplicación iterativa de la transformación.}  
    Sobre un clon del AST se aplica la política \emph{one-at-a-time} (véase \autoref{sec:one-at-a-time}), re-detectando tras cada cambio hasta alcanzar un punto fijo o el límite de seguridad.

    \item \textbf{Medición \emph{after}.}  
    Una vez finalizado el proceso iterativo, se vuelve a calcular la complejidad cognitiva y se obtiene $\Delta = \text{after} - \text{before}$.

    \item \textbf{Construcción del resultado experimental.}  
    Con la información anterior se construye una instancia inmutable de \texttt{ExperimentResult}.

    \item \textbf{Exportación de resultados.}  
    Finalmente, los resultados se exportan a JSON (incluyendo código fuente) y a CSV (formato tabular) (\autoref{fig:pipeline-det}).
\end{enumerate}

\begin{figure}[H]
\centering
\begin{tikzpicture}[
  term/.style={draw, rounded rectangle, fill=green!12, align=center,
               text width=3cm, inner sep=4pt, font=\small},
  endbad/.style={draw, rounded rectangle, fill=red!10, align=center,
                 text width=3cm, inner sep=4pt, font=\small},
  proc/.style={draw, fill=blue!5, align=center,
               text width=3.8cm, inner sep=5pt, font=\small},
  io/.style={draw, trapezium, trapezium left angle=70,
             trapezium right angle=110, trapezium stretches=true,
             fill=gray!15, align=center, text width=3.6cm, inner sep=4pt,
             font=\small},
  dec/.style={draw, diamond, aspect=2.2, fill=orange!12, align=center,
              text width=2.2cm, inner sep=1pt, font=\scriptsize},
  arr/.style={-{Latex}},
  lbl/.style={font=\scriptsize, inner sep=1.5pt},
]
% Eje principal
\node[term] (start) at (0,0)      {Inicio: \texttt{ExperimentCase}};
\node[dec]  (dpar)  at (0,-2.3)   {\textquestiondown{}AST válido y con método?};
\node[proc] (db)    at (0,-4.6)   {Detección con motivos\\+ CC \emph{before} (original)};
\node[proc] (loop)  at (0,-6.2)   {Detectar oportunidades\\(sobre el clon)};
\node[dec]  (dloop) at (0,-8.5)   {\textquestiondown{}oportunidad y pasadas $<$ \texttt{MAX}?};
\node[proc] (ar)    at (0,-10.8)  {CC \emph{after}, $\Delta$\\y \texttt{ExperimentResult}};
\node[io]   (exp)   at (0,-12.6)  {Exportar JSON + CSV};
\node[term] (fin)   at (0,-14.2)  {Fin};
% Ramas laterales
\node[endbad] (inel) at (4.8,-2.3) {Inelegible\\(\texttt{CC before = -1})};
\node[proc] (apply)  at (4.8,-8.5) {Aplicar la primera oportunidad};
% Flechas del eje
\draw[arr] (start) -- (dpar);
\draw[arr] (dpar)  -- node[lbl,left]{sí} (db);
\draw[arr] (db)    -- (loop);
\draw[arr] (loop)  -- (dloop);
\draw[arr] (dloop) -- node[lbl,left]{no} (ar);
\draw[arr] (ar)    -- (exp);
\draw[arr] (exp)   -- (fin);
% Ramas
\draw[arr] (dpar)  -- node[lbl,above]{no} (inel);
\draw[arr] (dloop) -- node[lbl,above]{sí} (apply);
\draw[arr] (apply) |- node[lbl,above]{re-detectar} (loop.east);
\end{tikzpicture}
\caption{Flujo de datos del pipeline determinista
(\texttt{BatchRunner.processCase}). Se hacen explícitos la rama de
inelegibilidad (fallo de parseo o ausencia de método $\rightarrow$
\texttt{CC before = -1}; el resultado inelegible también se registra y
exporta) y el bucle de la política \emph{one-at-a-time} (detectar
$\rightarrow$ aplicar la primera oportunidad $\rightarrow$ re-detectar,
hasta punto fijo o \texttt{MAX\_PASSES}). La CC \emph{before} se mide sobre
el método original y la \emph{after} sobre el clon transformado.}
\label{fig:pipeline-det}
\end{figure}

\subsection{Flujo del subsistema LLM (RQ3)}

Para RQ3 se añade un flujo paralelo coordinado por \texttt{CampaignExecutor} y \texttt{LlmCampaignRunner}, que se organiza en las siguientes etapas:

\begin{enumerate}
    \item \textbf{Carga del subset de evaluación.}  
    Se carga el subconjunto RQ3 definido en \texttt{LlmEvaluationSubset}, junto con sus correspondientes baselines deterministas.

    \item \textbf{Construcción del prompt.}  
    El mensaje enviado al modelo se genera mediante \texttt{LlmPromptBuilder}, utilizando la versión fijada del protocolo (\lstinline|PROMPT_VERSION = "v1.0"|).

    \item \textbf{Invocación al modelo.}  
    La ejecución del modelo se realiza a través de la interfaz \texttt{LlmResponseProvider}, con implementaciones específicas como \texttt{LiveLlmResponseProvider} y \texttt{PrerecordedResponseProvider}.

    \item \textbf{Validación mediante el oráculo.}  
    Cada respuesta se procesa con \texttt{LlmResponseValidator}, que aplica el oráculo experimental descrito en la \autoref{fig:oraculo}.

    \item \textbf{Agregación y exportación de resultados.}  
    Finalmente, los resultados se consolidan y exportan en artefactos.
\end{enumerate}

\subsection{Decisiones estructurales}
\begin{itemize}
  \item \textbf{Inmutabilidad de los resultados.} \texttt{ExperimentResult}
  y \texttt{LlmEvaluationResult} se construyen una sola vez y se exportan
  sin mutación posterior, lo que facilita la trazabilidad.
  \item \textbf{Provider intercambiable para LLM.}
  \texttt{LlmResponseProvider} permite ejecutar la misma campaña con
  respuestas reales o grabadas, desacoplando la lógica experimental del
  coste y la variabilidad de las llamadas remotas.
  \item \textbf{Exportadores múltiples.} CSV, JSON y Markdown cubren usos
  distintos (tabla agregada, evidencia completa, lectura humana).
  \item \textbf{Sin acoplamiento a SonarQube en el pipeline determinista.}
  El prototipo \emph{no invoca} SonarQube en el flujo principal; la
  validación cruzada se realiza de forma independiente del pipeline, como un
  contraste externo (\autoref{sec:sonar-val}).
\end{itemize}


\section{Implementación del prototipo}
\label{sec:implementacion}

Si la sección anterior describía el esqueleto, esta lo recubre de músculo.
Recorremos las piezas en el mismo orden en que intervienen en el flujo:
primero el detector, que decide qué se puede combinar; después el
transformador, que efectúa la combinación; luego el cálculo de la
complejidad cognitiva, con el que medimos el efecto; y, por último, el
pipeline experimental y el subsistema de evaluación de LLMs. 

\subsection{Detector}

La clase \texttt{NestedIfDetector} recorre el AST del método y visita todos
los \lstinline|if|, incluidos los anidados a cualquier profundidad. Para
cada \lstinline|if| candidato aplica en orden las precondiciones:

\begin{enumerate}
    \item \textbf{P1.} El \lstinline|if| externo no tiene rama \lstinline|else|.

    \item \textbf{P2--P3.} Su rama \lstinline|then| consiste en una única
    sentencia y esa sentencia es, a su vez, un \lstinline|if|. Se aceptan
    tanto el bloque con llaves de una sola sentencia como la forma sin
    llaves; en cualquier otro caso, el candidato se descarta por P2 o P3.

    \item \textbf{P4.} El \lstinline|if| interno tampoco tiene rama
    \lstinline|else|.

    \item \textbf{P5.} Ninguna de las dos condiciones contiene efectos
    colaterales: las cinco categorías descritas en \autoref{sec:precondiciones}.
\end{enumerate}

El detector ofrece dos modos de salida: uno devuelve solo las oportunidades
aceptadas y otro (\lstinline|detectWithReasons|) devuelve también los
candidatos rechazados con su motivo de descarte. Los diez motivos posibles
se agrupan en \emph{estructurales} (incumplimiento de P1--P4 o ausencia de
patrón combinable) y de \emph{efecto colateral} (las cinco variantes de P5);
el catálogo completo está en el Apéndice~D y su árbol de decisión en la
\autoref{fig:arbol-p1-p5}. Como los motivos se recogen por candidato y se
agregan a nivel de método, un mismo método puede aportar varias categorías
a \lstinline|discardCategories|, lo que constituye evidencia agregada para
RQ1.

\subsection{Transformador}
La clase \texttt{NestedIfTransformer} implementa el algoritmo de
\autoref{sec:diseno}. Algunos detalles relevantes son: \textbf{clonación}
para no dejar el
AST inconsistente; \textbf{parentizado selectivo},
solo si la expresión es una expresión binaria con operador
\lstinline|OR|; \textbf{normalización del cuerpo}, produciendo siempre código con llaves;
\textbf{eliminación defensiva del \texttt{else}};
y \textbf{verificación previa}, replicando
P1 y P4 como salvaguarda.

\subsection{Cálculo proxy de complejidad cognitiva}
\label{sec:proxy}
La clase \texttt{CognitiveComplexityCalculator} calcula la complejidad
cognitiva de un método siguiendo el modelo de SonarSource
\cite{10.1145/3194164.3194186}. Implementa las tres reglas del modelo:

\begin{itemize}
  \item \textbf{Incremento estructural (+1)}: \lstinline|if|,
  \lstinline|else if|, \lstinline|else|, \lstinline|for|, \emph{for-each},
  \lstinline|while|, \lstinline|do-while|, \lstinline|switch|,
  \lstinline|catch|, \lstinline|break|/\lstinline|continue| con etiqueta y
  el operador ternario \lstinline|?:|.
  \item \textbf{Incremento de anidamiento (+nivel)}: \lstinline|if|,
  \lstinline|for|, \emph{for-each}, \lstinline|while|, \lstinline|do-while|,
  \lstinline|switch| y \lstinline|catch| (no en \lstinline|else if| ni
  \lstinline|else|). Además, los cuerpos de \emph{lambdas}, clases anónimas
  y clases locales se procesan a nivel de anidamiento $+1$.
  \item \textbf{Operadores lógicos}: $+1$ por cada secuencia de operadores
  del mismo tipo (\lstinline|&&| o \lstinline=||=) en \emph{cualquier}
  expresión (condiciones, \lstinline|return|, asignaciones, argumentos,
  etc.). El conteo recorre la expresión de forma plana ignorando por
  completo los paréntesis, tal como especifica el modelo de referencia, de
  modo que las secuencias alternantes (p.\,ej. \lstinline=a || b && c || d=)
  se contabilizan correctamente.
\end{itemize}

\paragraph{Limitaciones.}
La única regla del modelo no implementada es la \textbf{detección de recursión} (que
SonarSource penaliza con $+1$). Por lo demás, el cálculo cubre el modelo
completo.

\paragraph{Validación frente a los casos oficiales de SonarSource.}
La
corrección del proxy se comprueba a nivel unitario contra los casos de
prueba \emph{oficiales} del modelo: se ha implementado una clase que compara, con
aserciones de igualdad exacta, el valor calculado frente al valor de
referencia de SonarSource para cada caso del \emph{fixture}
\texttt{CognitiveComplexityMethodCheck}, el mismo conjunto que
emplea la herramienta \texttt{SoftwareCognitiveComplexityReducer}. Los
casos cubren operadores lógicos (incluidas las secuencias alternantes con
paréntesis ignorados), estructuras de control (\lstinline|switch|,
\lstinline|try|/\lstinline|catch|/\lstinline|finally|,
\lstinline|break|/\lstinline|continue| etiquetados,
\lstinline|if|/\lstinline|else if|/\lstinline|else| anidados y bucles) y
las características antes limitadas (ternario, \emph{lambda}, clase anónima
y clase local). El proxy reproduce el valor oficial en todos los casos
cubiertos; la única exclusión es la recursión.

\subsection{Pipeline experimental}
El pipeline experimental, implementado por \texttt{BatchRunner}, recorre el
corpus caso por caso y aplica las oportunidades \emph{de una en una} (la
política iterativa de la \autoref{sec:one-at-a-time}). Para no alterar el método original, mide su
complejidad cognitiva \emph{antes}, lo clona y transforma el clon, sobre el
que mide la complejidad \emph{después}; la diferencia entre ambas es el
$\Delta$ del caso. Además de transformar, anota \emph{por qué} no combinó
los condicionales que descartó: esos motivos de rechazo se recogen aunque
el caso acabe siendo elegible, porque constituyen la evidencia de RQ1. Como
cada pasada aplica una sola oportunidad y vuelve a detectar, el número de
oportunidades detectadas coincide con el de aplicadas, y se registra
cuántas pasadas hicieron falta (con un tope de seguridad de diez). Si un
fragmento no se puede analizar —porque no compila o no contiene ningún
método— se marca como no elegible, se le asigna una complejidad
\emph{antes} de $-1$ como valor centinela y se anota el motivo.

\subsection{Subsistema de evaluación de LLMs (RQ3)}
\label{sec:llm-subsistema}
El paquete \texttt{es.tfm.refactoring.llm} implementa el protocolo de RQ3.
La \autoref{fig:secuencia-rq3} resume el ciclo completo de una evaluación;
los apartados siguientes detallan cada pieza (prompt, \emph{provider},
oráculo y campaña).

\begin{figure}[H]
\centering
\includegraphics[width=0.92\linewidth]{fig-secuencia-rq3}
\caption[Secuencia de una evaluación RQ3]{Diagrama de secuencia de una
evaluación RQ3: del caso y su \emph{baseline} (de RQ2) al veredicto, pasando
por prompt, \emph{provider}, modelo y oráculo.}
\label{fig:secuencia-rq3}
\end{figure}

\paragraph{Prompt.} \texttt{LlmPromptBuilder} contiene literalmente la versión \lstinline|v1.0| del prompt: un \emph{system prompt} con las reglas de
combinación (P1--P4 reescritas como restricciones), exigencia de
preservación de orden de evaluación y semántica, prohibición de modificar
firmas o código fuera de los \lstinline|if|, y opción de rechazo explícito
mediante \lstinline|NO_REFACTORING_APPLICABLE|. El \emph{user prompt}
envuelve el código del caso en un bloque \lstinline|```java|.

\paragraph{Protocolo.} El protocolo (\texttt{LlmExperimentProtocol})
encapsula los parámetros como configuración inmutable. Declara los modelos
\lstinline|gpt-4o| y \lstinline|gpt-4.1|; temperatura \lstinline|0.0|; 3
intentos por caso; prompt \lstinline|v1.0|; y el máximo de tokens de salida
en 2048. Un protocolo auxiliar mantiene los parámetros pero
restringe la lista a \lstinline|gpt-4o|, y 
\emph{no sustituye} al protocolo formal de dos modelos. Existe además un
tercer protocolo, idéntico en parámetros
salvo la versión del prompt (\lstinline|v2.0|, \emph{few-shot}); donde su
finalidad es una comparación exploratoria \emph{zero-shot} (v1.0) frente a
\emph{few-shot} (v2.0) pero \textbf{no} forma parte de la campaña de
referencia.

\paragraph{Provider.} La interfaz \texttt{LlmResponseProvider} admite dos
implementaciones intercambiables: la invocación real a la API de OpenAI
(con credencial) y las respuestas grabadas para \emph{dry runs} y tests.

\paragraph{Oráculo}
\label{sec:oraculo}

La validación de las respuestas generadas por los modelos
(\texttt{LlmResponseValidator}) se implementa como una cadena de
comprobaciones secuenciales:

\begin{enumerate}
    \item \textbf{Detección de rechazo explícito.}  
    Se comprueba si la respuesta del modelo contiene el marcador \lstinline|NO_REFACTORING_APPLICABLE|.

    \item \textbf{Extracción del bloque de código.}  
    Se extrae el bloque \lstinline|```java| mediante una expresión regular.

    \item \textbf{Comprobación de parseabilidad.}  
    El código extraído se analiza con JavaParser para verificar que sea sintácticamente válido.

    \item \textbf{Existencia de al menos un método.}  
    Se comprueba que el fragmento parseado contenga al menos un método sobre el que aplicar la validación.

    \item \textbf{Preservación de la firma.}  
    Se verifica que la firma del método coincida con la del caso original. Si difiere, se registra como error de validación.

    \item \textbf{Medición de complejidad cognitiva \emph{after}.}  
    Una vez aceptado el código desde el punto de vista estructural, se calcula la complejidad cognitiva resultante.

    \item \textbf{Cálculo del delta y comparación con el baseline.}  
    Se obtiene el valor $\Delta$ y se compara con el delta del baseline determinista.

    \item \textbf{Asignación del veredicto final.}  
    A partir de las comprobaciones anteriores se asigna el veredicto correspondiente.
\end{enumerate}

El conjunto de veredictos (\texttt{LlmVerdict}) está formado por \texttt{SUCCESS}, \texttt{INCORRECT}, \texttt{INVALID\_OUTPUT}, \texttt{REFUSED}, \texttt{PARTIAL} y \texttt{ERROR}.

El oráculo incorpora además la \textbf{guarda de elegibilidad v1.1}: si el baseline declara que un caso no es elegible (porque no se cumple alguna de las precondiciones P1--P5) y el LLM aplica una transformación con \lstinline|delta != 0|, el veredicto pasa a ser \texttt{INCORRECT}, aunque la complejidad disminuya. Con esta comprobación se evita clasificar como aciertos transformaciones que violan las precondiciones del detector.

La \autoref{fig:oraculo} resume esta cadena de validación en forma de diagrama de flujo, mientras que la \autoref{tab:oraculo} recoge las reglas exactas de asignación de veredicto.

\begin{figure}[H]
\centering
\begin{tikzpicture}[
  term/.style={draw, rounded rectangle, fill=green!12, align=center,
               text width=3.1cm, inner sep=4pt, font=\small},
  endv/.style={draw, rounded rectangle, fill=gray!12, align=center,
               text width=3.3cm, inner sep=4pt, font=\small},
  proc/.style={draw, fill=blue!5, align=center, text width=4.8cm,
               inner sep=5pt, font=\small},
  dec/.style={draw, diamond, aspect=2.4, fill=orange!12, align=center,
              text width=2.5cm, inner sep=1pt, font=\scriptsize},
  arr/.style={-{Latex}}, lbl/.style={font=\scriptsize, inner sep=1.5pt},
]
\node[term] (start) at (0,0)     {Respuesta del LLM};
\node[dec]  (d1)    at (0,-2.3)  {\textquestiondown{}rechazo explícito?};
\node[proc] (p1)    at (0,-4.6)  {Extraer bloque \lstinline|```java|};
\node[dec]  (d2)    at (0,-6.9)  {\textquestiondown{}bloque, parsea y $\geq 1$ método?};
\node[proc] (p2)    at (0,-9.4)  {Comprobar firma; medir CC \emph{after}; $\delta$ y comparación con baseline};
\node[term] (end)   at (0,-11.7) {Asignar veredicto (Tabla~\ref{tab:oraculo})};
% Salidas laterales
\node[endv] (refv)  at (5.4,-2.3) {Veredicto de rechazo (Tabla~\ref{tab:oraculo})};
\node[endv] (inv)   at (5.4,-6.9) {\texttt{INVALID\_OUTPUT}};
% Flechas
\draw[arr] (start) -- (d1);
\draw[arr] (d1) -- node[lbl,left]{no} (p1);
\draw[arr] (p1) -- (d2);
\draw[arr] (d2) -- node[lbl,left]{sí} (p2);
\draw[arr] (p2) -- (end);
\draw[arr] (d1) -- node[lbl,above]{sí} (refv);
\draw[arr] (d2) -- node[lbl,above]{no} (inv);
\end{tikzpicture}
\caption{Cadena de validación del oráculo \texttt{LlmResponseValidator}
(diagrama de flujo). Las salidas tempranas cubren el rechazo explícito
(\lstinline|NO_REFACTORING_APPLICABLE|) y la salida no válida
(\texttt{INVALID\_OUTPUT}); el resto desemboca en la asignación de veredicto
según la \autoref{tab:oraculo}.}
\label{fig:oraculo}
\end{figure}

\begin{table}[H]
\centering
\caption{Reglas de asignación de veredicto del oráculo (política v1.1).
Fuente: \texttt{LlmResponseValidator.evaluate}. $\delta$ es la variación de
CC del LLM y \emph{baseline} la del prototipo determinista.}
\label{tab:oraculo}
\begin{tabular}{@{}p{0.46\linewidth}p{0.30\linewidth}l@{}}
\toprule
\textbf{Situación} & \textbf{Condición} & \textbf{Veredicto} \\
\midrule
Rechazo explícito, caso inelegible & rechazo $\wedge$ \lstinline|!eligible| & \texttt{SUCCESS} \\
Rechazo explícito, caso elegible & rechazo $\wedge$ \lstinline|eligible| & \texttt{REFUSED} \\
Sin bloque, no parsea o sin método & --- & \texttt{INVALID\_OUTPUT} \\
Firma modificada u otro error & errores $\neq \emptyset$ & \texttt{INCORRECT} \\
Inelegible y transforma (guarda v1.1) & \lstinline|!eligible| $\wedge\ \delta \neq 0$ & \texttt{INCORRECT} \\
Inelegible y no transforma & \lstinline|!eligible| $\wedge\ \delta = 0$ & \texttt{SUCCESS} \\
Elegible, $\delta<0$ igual o mayor que baseline & \lstinline|eligible| $\wedge\ \delta \le$ base $< 0$ & \texttt{SUCCESS} \\
Elegible, reduce menos que baseline & \lstinline|eligible| $\wedge$ base $< \delta < 0$ & \texttt{PARTIAL} \\
Elegible, no reduce ($\delta=0$, base $<0$) & \lstinline|eligible| $\wedge\ \delta = 0$ & \texttt{INCORRECT} \\
Error técnico (red, \emph{timeout}\dots) & excepción en el \emph{provider} & \texttt{ERROR} \\
\bottomrule
\end{tabular}
\end{table}

\paragraph{Campaña y exportación.} La campaña (\texttt{LlmCampaignRunner})
orquesta la ejecución sobre el subset, genera un registro por llamada
(timestamp, versión de modelo, prompt enviado, tiempo de respuesta) y delega
en el validador; las excepciones del \emph{provider} se propagan como
veredicto \texttt{ERROR}. El punto de
entrada ofrece \lstinline|--mode=live| (provider
real, validación de credenciales obligatoria) y \lstinline|--mode=dry-run|
(provider grabado, sin red).


\chapter{Diseño experimental}
\label{sec:experimento}

Este capítulo detalla cómo se han diseñado y ejecutado los experimentos que
responden a RQ2 y RQ3. En primer lugar se describe el corpus sobre el que se
trabaja. A continuación se definen las métricas y la línea base con las que se
mide el efecto de la transformación. Después se presentan, por separado, el
protocolo determinista de RQ2 y el protocolo de evaluación de los modelos de
lenguaje de RQ3, y el capítulo se cierra con la validación complementaria de la
métrica frente a SonarQube. Un principio recorre todo el capítulo: la
\textbf{reproducibilidad}. Cada decisión experimental se fija de antemano y se
registra, de modo que el procedimiento aquí descrito pueda repetirse y conducir
al mismo resultado.

\section{Corpus}
\label{sec:corpus}
El corpus se carga, por defecto, desde recursos empaquetados mediante
\texttt{PilotCorpusLoader} y \texttt{RealDatasetLoader}; ambos cargadores
admiten además un \emph{directorio del sistema de archivos} ---con subcarpetas
\texttt{pilot-corpus/} y \texttt{real-corpus/}---, de modo que el corpus puede
sustituirse, sin recompilar, indicando esa ruta al ejecutar el lote. Su
composición se resume en la \autoref{tab:corpus}.

\begin{table}[ht]
\centering
\caption{Composición del corpus de los experimentos RQ1/RQ2.}
\label{tab:corpus}
\begin{tabular}{@{}lrl@{}}
\toprule
\textbf{Subconjunto} & \textbf{Casos} & \textbf{Finalidad} \\
\midrule
Piloto (sintético) & 8   & Casos canónicos y un rechazo por cada precondición \\
Real (código abierto) & 118 & Validación sobre código real de proyectos de código abierto \\
\midrule
\textbf{Total} & \textbf{126} & 111 elegibles, 15 no elegibles \\
\bottomrule
\end{tabular}
\end{table}

\paragraph{Corpus piloto}:

Ficheros sintéticos que cubren los casos canónicos del patrón objetivo y
sus rechazos típicos (casos válidos, inválidos por violación de cada
precondición y un caso mixto con varias oportunidades). Su finalidad es
servir de prueba controlada del pipeline.

\paragraph{Corpus real}:

El corpus real está formado por fragmentos extraídos de proyectos Java de código abierto. Cada caso conserva metadatos que permiten trazar su procedencia y contexto original. No son anotaciones de Java ni proceden de ninguna biblioteca: son etiquetas propias del proyecto, escritas como comentarios de línea al inicio de cada fichero (\lstinline|@project|, \lstinline|@file|, \lstinline|@method|, \lstinline|@license| y \lstinline|@sonarCCBefore|), que el cargador del corpus lee y asocia a cada caso.

Los proyectos fuente proceden de dos orígenes principales:

\begin{enumerate}
    \item \textbf{Proyectos de código abierto de uso general}, añadidos en este trabajo para ampliar la
    variedad de dominios y estilos de código. Aportaron casos Apache Ant,
    Apache POI, Jackson Core y Jackson Databind (FasterXML) y las bibliotecas
    Apache Commons (\emph{Lang}, \emph{Collections}, \emph{Math} e \emph{IO}).
    Otros proyectos de este grupo se escanearon pero no aportaron ningún caso
    (Apache Commons \emph{Numbers}, \emph{Geometry} y \emph{Compress}, ASM y
    PDFBox). El escaneo se realizó en varias tandas.

    \item \textbf{Proyectos de la evaluación de Saborido et
    al.}~\cite{10.1109/ACCESS.2022.3144743}: los diez proyectos del antecedente.
    Aportaron casos \texttt{cybercaptor-server}, \texttt{MOEAFramework},
    \texttt{knowage-server}, \texttt{fastjson}, \texttt{iotbroker}, jMetal y
    Bytecode Viewer; en cambio, \texttt{fiware-commons}, \texttt{jedis} y
    \texttt{aioteslake} se escanearon pero no aportaron ninguno.
\end{enumerate}

El detalle por proyecto (dominio, volumen de código analizado y casos
elegibles y no elegibles) se presenta junto al análisis de RQ1
(\autoref{sec:resultados-rq1}).

El corpus real no procede de un
muestreo arbitrario, sino del escaneo sistemático de los proyectos citados. En
conjunto, el escaneo evaluó 82\,655 candidatos a \texttt{if} externo del patrón
en 23 proyectos de código abierto y localizó 117 instancias elegibles del patrón
\texttt{if}--\texttt{if} combinable. A partir de esas instancias (junto con
casos no elegibles representativos y los 8 casos del piloto sintético) se
consolidó el corpus de 126 casos, de los cuales 99 son casos reales elegibles.
El corpus es, por tanto, esencialmente un censo de las oportunidades del patrón
presentes en los proyectos analizados, y no una submuestra de los
\emph{~}1\,000 métodos con $\text{CC}>15$ considerados en el anteproyecto.

\paragraph{Nota de alcance respecto al anteproyecto.} El anteproyecto aprobado
preveía partir de los \emph{~}1\,000 métodos con $\text{CC}>15$ del repositorio
de Saborido et al., sin acotar el tipo de transformación. Al concretar el
alcance en la combinación de condicionales anidados, un patrón específico y
verificable, no la reducción de CC en general, el corpus resultante quedó en
126 casos: los métodos que contienen exactamente este patrón en los proyectos
fuente, más los 8 casos sintéticos del piloto. El resto de los métodos con
$\text{CC}>15$ de esos proyectos pertenece a familias de transformación
distintas (principalmente extracción de método) que quedan fuera del alcance
declarado en \autoref{sec:alcance}.


\paragraph{Cobertura de RQ3.} La evaluación de RQ3 se ejecuta sobre el
\textbf{corpus consolidado completo} (126 casos: 8 piloto y 118 reales; 104
elegibles y 22 no elegibles), de modo que cubre todos los ejes de variación
(casos elegibles, parciales y no elegibles) y los principales motivos de descarte estructurales.
La \autoref{tab:subset-rq3} recoge una selección representativa de casos.

\begin{table}[ht]
\centering
\footnotesize
\caption{Ejemplos representativos de los casos de evaluación de RQ3
(la campaña cubre el corpus completo de 126 casos).}
\label{tab:subset-rq3}
\begin{tabularx}{\linewidth}{@{}lll>{\raggedright\arraybackslash}X@{}}
\toprule
\textbf{Identificador} & \textbf{Origen} & \textbf{Categoría} & \textbf{Justificación} \\
\midrule
\texttt{PILOT\_VALID\_SIMPLE} & Piloto & Elegible & Caso base \lstinline|if + if|. \\
\texttt{PILOT\_VALID\_NESTED\_IN\_LOOP} & Piloto & Elegible & Anidamiento en \lstinline|for|. \\
\texttt{REAL\_COMMONS\_MATH\_CONVERGED} & Real & Elegible & Doble condición numérica. \\
\texttt{REAL\_ANT\_MATCH\_PATH} & Real & Elegible & Triple \lstinline|if|; se combinan los dos pares anidados. \\
\texttt{REAL\_COMMONS\_MATH\_VALIDATE\_RANGE} & Real & No elegible (P1) & \lstinline|outer if| con \lstinline|else|. \\
\texttt{REAL\_COMMONS\_COLLECTIONS\_GET} & Real & Elegible & Llamada a método en la condición (combinable). \\
\bottomrule
\end{tabularx}
\end{table}

\section{Métricas y línea base}
La evaluación se apoya en dos métricas de complejidad y una referencia de
comparación, todas fijadas de antemano para garantizar la reproducibilidad:
\begin{itemize}
  \item \textbf{Métrica primaria:} complejidad cognitiva proxy (CC) calculada
  por \texttt{Cognitive\allowbreak Complexity\allowbreak Calculator} antes y después de la
  transformación; se reporta \emph{before}, \emph{after} y
  $\Delta = \text{after} - \text{before}$ (un $\Delta$ negativo indica
  reducción). Las limitaciones del proxy están en \autoref{sec:proxy}.
  \item \textbf{Métrica complementaria:} complejidad cognitiva medida por
  SonarQube Community Build sobre el corpus completo, que sirve para validar el
  proxy (\autoref{sec:sonar-val}).
  \item \textbf{Baseline determinista:} para cada caso, el resultado del
  pipeline (\texttt{BatchRunner}) constituye el baseline contra el que se
  contrastan las salidas de los LLMs en RQ3; declara explícitamente la
  elegibilidad y, en su caso, el $\Delta$ aplicable.
\end{itemize}

\section{Validación del proxy frente a SonarQube}
\label{sec:sonar-val}
Como la métrica primaria es un \emph{proxy} y no la implementación oficial,
conviene comprobar su validez antes de emplearla en los experimentos. Para ello
se contrastan sus valores con los de SonarQube Community Build (versión
\texttt{26.4.0.121862}) sobre el \textbf{corpus completo}.

El contraste se realiza caso por caso: se materializan los métodos \emph{before}
y \emph{after} como ficheros Java independientes, se analizan en dos proyectos
de SonarQube y se recupera su complejidad cognitiva oficial por fichero mediante
la API, que se compara con la del proxy.

El procedimiento es reproducible desde PowerShell: el análisis se lanza con el
\emph{goal} \texttt{mvn sonar:sonar} contra un servidor local de SonarQube,
arrancado en un contenedor \emph{Docker} mediante \texttt{docker-compose}. Los
resultados de este contraste se presentan en la \autoref{sec:resultados-sonar}.

\section{Protocolo RQ2 — Batch determinista \emph{before/after}}
\label{sec:exp-rq2}
El protocolo consta de cuatro pasos:

\begin{enumerate}
    \item Cargar la totalidad del corpus piloto y real.
    \item Para cada \texttt{ExperimentCase}, ejecutar
    \lstinline|BatchRunner.processCase| con el detector en modo
    \texttt{STRUCTURAL} (precondiciones estructurales P1--P4) y la política
    una por una.
    \item Recoger métricas \emph{before}, \emph{after} y $\Delta$, y motivos de
    descarte para los no elegibles.
    \item Exportar a JSON y CSV.
\end{enumerate}

De este protocolo se \textbf{congelan} ---es decir, se fijan de antemano y no
varían entre ejecuciones--- cuatro elementos: la composición del corpus, el modo
de detección (\texttt{STRUCTURAL}), el conjunto de precondiciones P1--P4 y la
política de aplicación una por una. Cualquier cambio en
\texttt{Cognitive\allowbreak Complexity\allowbreak Calculator} invalidaría las métricas y exigiría
repetir la ejecución completa.

\section{Protocolo RQ3 — Evaluación de LLMs}
\label{sec:exp-rq3}
El protocolo está congelado en
\lstinline|LlmExperimentProtocol.defaultProtocol()|
(\autoref{tab:protocolo-rq3}).

\begin{table}[ht]
\centering
\caption{Protocolo congelado de RQ3.}
\label{tab:protocolo-rq3}
\begin{tabular}{@{}lll@{}}
\toprule
\textbf{Parámetro} & \textbf{Valor} & \textbf{Justificación} \\
\midrule
Modelos & \texttt{gpt-4o}, \texttt{gpt-4.1} & Comparativa intra-proveedor (OpenAI). \\
Temperatura & \texttt{0.0} & Determinismo esperado. \\
Intentos por caso & \texttt{3} & Detección de inconsistencia residual. \\
Prompt & \texttt{v1.0} & Trazabilidad (texto literal en código). \\
\emph{Max tokens} salida & \texttt{2048} & Métodos refactorizados completos. \\
Cobertura & 126 casos & Corpus completo (elegibles/parciales/no elegibles). \\
Total invocaciones & 756 & $126 \times 2 \times 3$. \\
\bottomrule
\end{tabular}
\end{table}

La comparación se restringe deliberadamente a un único proveedor (OpenAI) por
tres razones. Primera, de \emph{validez interna}: contrastar dos modelos del
mismo proveedor (\texttt{gpt-4o} y \texttt{gpt-4.1}) mantiene constantes la
API, el tokenizador y la vía de acceso, de modo que las diferencias observadas
se atribuyen al modelo y no al entorno. Segunda, de \emph{alcance y recursos}:
el trabajo dispone de una única clave de API financiada y de un presupuesto
temporal acotado, lo que hace inviable una campaña multiproveedor sin
comprometer la profundidad del resto del estudio. Tercera, de
\emph{representatividad}: los modelos elegidos pertenecen a la familia de
propósito general de referencia en generación de código (linaje de Codex,
\autoref{sec:llm-antecedentes}). La generalización a otras familias de modelos
queda explícitamente como amenaza a la validez externa y trabajo futuro
(\autoref{cap:amenazas}); el diseño lo facilita, pues el \emph{provider} del
protocolo es intercambiable y añadir otro proveedor es un cambio de
configuración, no de arquitectura.

El procedimiento que se aplica a cada caso consta de los siguientes pasos:
\begin{enumerate}
  \item cargar el caso y su \emph{baseline} determinista;
  \item construir el \emph{prompt} \texttt{v1.0};
  \item invocar el modelo;
  \item almacenar la respuesta cruda;
  \item aplicar el oráculo (\autoref{fig:oraculo});
  \item registrar el \texttt{LlmEvaluationResult}.
\end{enumerate}
Estos pasos se repiten para cada intento y cada modelo. Como en RQ2, el
protocolo está \textbf{congelado}: se fijan de antemano los modelos, la
temperatura, el número de intentos, la versión del \emph{prompt}, el oráculo y
la composición del conjunto evaluado. El único componente intercambiable es el
\emph{proveedor} de respuestas, cuyo modo oficial es \lstinline|live| (la
invocación real a la API); esa indirección permite, por ejemplo, sustituir las
llamadas por respuestas grabadas en las pruebas sin alterar el resto del
protocolo.

\section{Reproducibilidad}
La reproducibilidad del estudio no se enuncia como aspiración, sino que se apoya
en elementos concretos y verificables:
\begin{itemize}
  \item exigencia explícita de \lstinline|OPENAI_API_KEY| en modo \lstinline|live|;
  \item versionado del \emph{prompt} propagado a cada artefacto;
  \item versiones fijadas en \texttt{pom.xml} (Java 17, JavaParser 3.26.4,
  JUnit 5.11.4, Gson 2.11.0);
  \item artefactos persistentes auditables sin ejecutar el código.
\end{itemize}


\section{Resultados}
\label{sec:resultados}

Este capítulo presenta los resultados siguiendo el orden de las preguntas de
investigación: primero qué casos resultaron combinables y por qué se
descartaron los demás (RQ1); después con qué intensidad bajó la complejidad
cognitiva en los casos elegibles (RQ2), con su validación de contraste frente
a SonarQube; y por último cómo se comportaron los modelos de lenguaje sobre el
mismo problema (RQ3). Todas las cifras de complejidad proceden del proxy del
prototipo (\autoref{sec:proxy}) y de los artefactos generados; los estadísticos descriptivos se han calculado sobre esos
mismos ficheros y ninguno se ha estimado a mano para esta memoria.

\subsection{RQ1 — Taxonomía de aplicabilidad}
\label{sec:resultados-rq1}
Sobre el corpus consolidado de 126 casos (8 piloto y 118 reales), el detector
marcó como elegibles 104 y descartó 22. La cifra interesa menos por sí misma
que por su lectura: en un corpus dominado por código real de proyectos
abiertos, una mayoría amplia de los patrones \lstinline|if|--\lstinline|if|
examinados sí admite la combinación bajo P1--P5.

Más revelador es el reparto de los motivos de descarte, que es justamente lo
que da respuesta empírica a RQ1: \emph{\textquestiondown{}en qué casos no se
pueden combinar?} La \autoref{tab:descartes} y la \autoref{fig:descartes}
reproducen la frecuencia de cada categoría tal como la emite el pipeline.
Conviene leerlas con una advertencia: los motivos se agregan por candidato y
un mismo método puede aportar varios, de modo que la suma supera el número de
casos inelegibles; no es un error, sino el modo en que
\lstinline|detectWithReasons| reporta la evidencia.

\begin{table}[ht]
\centering
\caption{Frecuencia de motivos de descarte en el corpus consolidado.
Fuente: \texttt{output/rq2-batch/rq2-summary.md}.}
\label{tab:descartes}
\begin{tabular}{@{}lr@{}}
\toprule
\textbf{Categoría (\texttt{DiscardReason})} & \textbf{Apariciones} \\
\midrule
\texttt{SINGLE\_STATEMENT\_NOT\_IF}            & 16 \\
\texttt{METHOD\_CALL\_IN\_CONDITION} (P5)      & 5 \\
\texttt{OUTER\_BLOCK\_MULTIPLE\_STATEMENTS} (P2) & 5 \\
(sin patrón de \lstinline|if| anidado)         & 5 \\
\texttt{OUTER\_HAS\_ELSE} (P1)                 & 4 \\
\texttt{INNER\_HAS\_ELSE} (P4)                 & 1 \\
\texttt{ASSIGNMENT\_IN\_CONDITION} (P5)        & 1 \\
\texttt{INCREMENT\_OR\_DECREMENT\_IN\_CONDITION} (P5) & 1 \\
\bottomrule
\end{tabular}
\end{table}

\begin{figure}[ht]
\centering
\begin{tikzpicture}
\begin{axis}[
  width=0.9\linewidth, height=5.8cm,
  xbar, xmin=0, xmax=19,
  xlabel={Apariciones},
  symbolic y coords={{Incr./decr. (P5)},{Asignación (P5)},{Else interno (P4)},{Else externo (P1)},{Sin patrón anidado},{Bloque múltiple (P2)},{Llamada a método (P5)},{Sentencia no-if}},
  ytick=data, nodes near coords, nodes near coords align={horizontal},
  enlarge y limits=0.12, y tick label style={font=\small},
]
\addplot+[fill=blue!25] coordinates {
  (1,{Incr./decr. (P5)}) (1,{Asignación (P5)}) (1,{Else interno (P4)})
  (4,{Else externo (P1)}) (5,{Sin patrón anidado}) (5,{Bloque múltiple (P2)})
  (5,{Llamada a método (P5)}) (16,{Sentencia no-if}) };
\end{axis}
\end{tikzpicture}
\caption{Distribución de los motivos de descarte (RQ1). Fuente:
\texttt{output/rq2-batch/rq2-summary.md}.}
\label{fig:descartes}
\end{figure}

De aquí se extraen dos lecturas. La primera es que el grueso de los descartes
no se debe a que la combinación sea insegura, sino a que el patrón objetivo
simplemente no está presente: la sentencia única del bloque externo no es un
\lstinline|if| anidado (\texttt{SINGLE\_STATEMENT\_NOT\_IF}, el motivo más
frecuente, con 16 apariciones). La segunda es que, cuando el patrón sí aparece
pero se rechaza, el bloqueo proviene sobre todo de P5 —y dentro de P5, de las
llamadas a método en la condición—, seguido de los bloques externos con más de
una sentencia. En otras palabras, en código real el principal enemigo de la
combinación segura no es el \lstinline|else|, sino el posible efecto colateral
escondido en la condición. Este resultado respalda, a posteriori, la decisión
de diseño de tratar P5 de forma sintáctica y conservadora.

\subsection{RQ2 — Impacto en la complejidad cognitiva}
\label{sec:resultados-rq2}
La respuesta cuantitativa a RQ2 es directa: sobre los 104 casos elegibles, la
complejidad cognitiva agregada (proxy) se reduce en $\Delta = -278$, lo que
supone una reducción relativa del 9{,}8\,\% sobre la complejidad total de esos
métodos (de 2825 a 2547 puntos). La reducción no se concentra en unos pocos
casos: los 104 elegibles presentan, sin excepción, un $\Delta$ negativo
(media $-2{,}67$ puntos por caso, desviación típica $2{,}29$). Los 22 casos
inelegibles mantienen, por construcción, $\Delta = 0$.

La \autoref{tab:rq2-cinco} resume la distribución completa mediante los
estadísticos descriptivos de la complejidad antes y después de refactorizar y
de su diferencia. La \autoref{fig:boxplots} los representa como diagramas de
caja y la \autoref{fig:hist-delta} muestra el histograma de las reducciones.

\begin{table}[ht]
\centering
\caption{Estadísticos descriptivos sobre los 104 casos elegibles (CC proxy):
mínimo, cuartiles, mediana, máximo, media y desviación típica. Cuartiles por
interpolación lineal. Fuente: \texttt{output/rq2-batch/rq2-all-results.csv}.}
\label{tab:rq2-cinco}
\begin{tabular}{@{}lrrrrrrrr@{}}
\toprule
\textbf{Métrica} & \textbf{n} & \textbf{mín} & \textbf{Q1} & \textbf{mediana} & \textbf{Q3} & \textbf{máx} & \textbf{media} & \textbf{sd} \\
\midrule
CC antes   & 104 & 3   & 4    & 10{,}5 & 20{,}25 & 228 & 27{,}16 & 46{,}01 \\
CC después & 104 & 2   & 3    & 8      & 18{,}25 & 220 & 24{,}49 & 44{,}44 \\
$\Delta$   & 104 & $-14$ & $-3$ & $-2$ & $-1$    & $-1$ & $-2{,}67$ & 2{,}29 \\
\bottomrule
\end{tabular}
\end{table}

El rango intercuartílico pasa de $16{,}25$ (antes) a $15{,}25$ (después): la
caja de la complejidad se desplaza hacia abajo tras refactorizar (la mediana
baja de $10{,}5$ a $8$ y el tercer cuartil de $20{,}25$ a $18{,}25$), con una
cola larga de métodos muy complejos —hasta 228, varios de Knowage— que
aparecen como valores atípicos según el criterio de Tukey. El histograma de
$\Delta$ confirma la asimetría del efecto: el 50\,\% central de las reducciones
se sitúa entre $-1$ y $-3$ (mediana $-2$), y la gran mayoría de los casos
ahorran uno o dos puntos, con una cola de reducciones mayores que llega a
$-14$.

\begin{figure}[ht]
\centering
\begin{tikzpicture}
\begin{axis}[
  width=0.52\linewidth, height=6cm,
  ymode=log, ymin=1, ymax=320,
  ylabel={Complejidad cognitiva (proxy, log)},
  xtick={1,2}, xticklabels={Antes, Después},
  boxplot/draw direction=y, log ticks with fixed point,
]
\addplot+[boxplot prepared={lower whisker=3, lower quartile=4, median=10.5, upper quartile=20.25, upper whisker=42}]
  coordinates {(0,46)(0,52)(0,53)(0,57)(0,63)(0,63)(0,157)(0,165)(0,166)(0,166)(0,167)(0,183)(0,183)(0,228)};
\addplot+[boxplot prepared={lower whisker=2, lower quartile=3, median=8, upper quartile=18.25, upper whisker=39}]
  coordinates {(0,44)(0,48)(0,52)(0,57)(0,58)(0,150)(0,158)(0,159)(0,159)(0,160)(0,176)(0,176)(0,220)};
\end{axis}
\end{tikzpicture}\hfill
\begin{tikzpicture}
\begin{axis}[
  width=0.40\linewidth, height=6cm,
  ylabel={$\Delta$ CC (proxy)}, xtick={1}, xticklabels={$\Delta$},
  boxplot/draw direction=y,
]
\addplot+[boxplot prepared={lower whisker=-6, lower quartile=-3, median=-2, upper quartile=-1, upper whisker=-1}]
  coordinates {(0,-14)(0,-9)(0,-8)(0,-8)(0,-8)(0,-7)(0,-7)(0,-7)(0,-7)(0,-7)(0,-7)(0,-7)};
\end{axis}
\end{tikzpicture}
\caption{Diagramas de caja (104 casos elegibles, CC proxy). Izquierda:
complejidad antes vs.\ después (escala logarítmica; los puntos son valores
atípicos según el criterio de Tukey de $1{,}5\cdot\mathrm{IQR}$). Derecha:
distribución de la reducción $\Delta$. Datos: \texttt{output/rq2-batch/}.}
\label{fig:boxplots}
\end{figure}

\begin{figure}[ht]
\centering
\begin{tikzpicture}
\begin{axis}[
  ybar, width=0.9\linewidth, height=5.5cm,
  xlabel={$\Delta$ de complejidad cognitiva}, ylabel={Nº de casos},
  xtick={-14,-12,-10,-8,-6,-4,-2}, ymin=0, ymax=40, bar width=7pt,
  nodes near coords, every node near coord/.append style={font=\scriptsize},
  enlarge x limits=0.04,
]
\addplot+[fill=blue!25] coordinates {
  (-14,1)(-13,0)(-12,0)(-11,0)(-10,0)(-9,1)(-8,3)(-7,7)(-6,1)(-5,3)(-4,5)(-3,12)(-2,34)(-1,37)};
\end{axis}
\end{tikzpicture}
\caption{Histograma de la reducción $\Delta$ sobre los 104 casos elegibles
(CC proxy). La masa se concentra en reducciones pequeñas ($-1$, $-2$), con una
cola de reducciones mayores. Datos: \texttt{output/rq2-batch/rq2-all-results.csv}.}
\label{fig:hist-delta}
\end{figure}

Más allá de la distribución, el detalle por caso
(\autoref{tab:rq2-destacados}) muestra dos comportamientos complementarios. En
métodos pequeños, combinar un único anidamiento ahorra uno o dos puntos. En
métodos grandes y muy anidados —habituales en el corpus real— el ahorro
absoluto crece, aunque sigue acotado por el hecho de que cada aplicación
elimina un solo nivel de anidamiento. El caso de mayor reducción individual es
\texttt{CYBERCAPTOR\_VERTEX\_GET\_RELATED\_MACHINE}, que baja de 52 a 38
($\Delta=-14$).

\begin{table}[ht]
\centering
\caption{Casos destacados de RQ2 (proxy). El conjunto completo de los 126
casos está en \texttt{output/rq2-batch/rq2-summary.md} y en el Apéndice~E.}
\label{tab:rq2-destacados}
\begin{tabular}{@{}lrrrr@{}}
\toprule
\textbf{caseId} & \textbf{Aplic.} & \textbf{CC antes} & \textbf{CC después} & \textbf{$\Delta$} \\
\midrule
\texttt{CYBERCAPTOR\_VERTEX\_GET\_RELATED\_MACHINE} & 1 & 52 & 38 & $-14$ \\
\texttt{ANT\_STRIP\_JAVA\_COMMENTS\_READ}           & 1 & 57 & 48 & $-9$ \\
\texttt{KNOWAGE\_DOCUMENT\_EXECUTION\_RESOURCE\_GET\dots} & 1 & 228 & 220 & $-8$ \\
\texttt{JACKSON\_CORE\dots SKIP\_COLON\_FAST}       & 2 & 36 & 28 & $-8$ \\
\texttt{REAL\_COMMONS\_LANG\_CONTAINS\_NONE}        & 1 & 15 & 11 & $-4$ \\
\texttt{PILOT\_VALID\_NESTED\_IN\_LOOP}             & 1 & 6  & 4  & $-2$ \\
\texttt{PILOT\_VALID\_SIMPLE}                       & 1 & 3  & 2  & $-1$ \\
\texttt{REAL\_COMMONS\_MATH\_VALIDATE\_RANGE} (P1)  & 0 & 4  & 4  & $0$ \\
\bottomrule
\end{tabular}
\end{table}

De esta observación se desprende un matiz honesto que conviene no maquillar: la
transformación mejora la legibilidad local, pero no \emph{arregla} un método
intrínsecamente complejo. Un método con complejidad 228 sigue siendo difícil
después de bajar a 220. El valor de la refactorización está en eliminar
anidamiento accidental allí donde lo hay, no en rescatar diseños que arrastran
problemas más profundos.

\subsection{Eliminación de issues de SonarQube}
\label{sec:resultados-issues}
La reducción de complejidad tiene una consecuencia práctica que va más allá del
$\Delta$. SonarQube marca un \emph{issue} de mantenibilidad cuando la
complejidad cognitiva de un método supera el umbral por defecto de 15 de su
regla S3776. Tiene sentido, por tanto, preguntarse en cuántos casos la
refactorización hace \emph{desaparecer} ese issue, es decir, cuántos métodos
cruzan de $\text{CC}>15$ a $\text{CC}\le 15$ (\autoref{tab:issues-sonar}).

\begin{table}[ht]
\centering
\caption{Cruce del umbral de issue de SonarQube ($\text{CC}=15$) en los 104
casos elegibles (CC proxy). Fuente: \texttt{output/rq2-batch/rq2-all-results.csv}.}
\label{tab:issues-sonar}
\begin{tabular}{@{}lr@{}}
\toprule
\textbf{Situación} & \textbf{Casos} \\
\midrule
Con issue antes ($\text{CC}>15$)                 & 41 \\
\quad $\rightarrow$ issue \textbf{eliminado} ($\text{CC}\le 15$ después) & 6 \\
\quad $\rightarrow$ issue persiste ($\text{CC}>15$ después)             & 35 \\
Sin issue ya antes ($\text{CC}\le 15$)           & 63 \\
\bottomrule
\end{tabular}
\end{table}

De los 104 casos elegibles, 41 superaban el umbral antes de refactorizar. Tras
aplicar la transformación, 6 de ellos bajan a 15 o menos —el issue desaparece
sin tocar la semántica— y 35 siguen por encima. Los 63 restantes ya estaban en
zona sin issue. Dicho de otro modo, la refactorización elimina el issue de
SonarQube en el 14{,}6\,\% de los métodos que lo tenían. La lectura honesta de
este dato encaja con la de RQ2: la combinación elimina el issue sobre todo
cuando el método estaba \emph{cerca} del umbral; en los métodos muy complejos
(los de Knowage, con CC entre 150 y 228) reduce la complejidad pero una sola
combinación no basta para cruzar el umbral. El efecto, por tanto, es local y
acotado, pero tangible: seis issues de mantenibilidad menos a coste cero en
términos de comportamiento.

\subsection{Validación de contraste con SonarQube}
\label{sec:resultados-sonar}
La métrica de RQ2 es un proxy, y por eso se contrasta con la implementación
oficial de SonarQube sobre un subset de cuatro casos (\autoref{sec:sonar-val}).
El artefacto \texttt{output/rq2-sonar-validation/rq2-sonar-validation.md}
recoge una coincidencia exacta entre proxy y SonarQube en las ocho medidas del
subset (cuatro \emph{before} y cuatro \emph{after}), con $\Sigma\Delta = -8$ en
ambos.

\pendienteverif{} \textbf{Este resultado no se da por confirmado.} Los
artefactos del repositorio son contradictorios: mientras la tabla anterior y la
cabecera del documento declaran una ejecución real completada, el fichero
\texttt{metadata.json} anota expresamente que \emph{«no se han ejecutado
escaneos reales de SonarQube en esta entrega»}, y la sección de trazabilidad
del propio informe deja la versión exacta, las fechas y el \emph{commit hash}
marcados como pendientes. Hasta que se ejecute la corrida real y se completen
esos campos con valores independientes del proxy, esta validación se presenta
como \textbf{provisional}. La incertidumbre afecta solo a este contraste; los
resultados de RQ2 (\autoref{sec:resultados-rq2}) se sostienen sobre el proxy y
son independientes de él.

\subsection{RQ3 — Refactorización por modelos de lenguaje}
\label{sec:resultados-rq3}
La campaña real (fase~9, modo \emph{live}) ejecutó las 300 invocaciones
previstas —50 casos $\times$ 2 modelos $\times$ 3 intentos— sin ningún
incidente técnico. El subconjunto de 50 casos (36 elegibles y 14 no elegibles)
amplía sustancialmente la evaluación respecto a versiones anteriores del
protocolo y permite analizar el comportamiento de los modelos no solo de forma
agregada, sino separando los casos que \emph{deben} transformarse de los que
\emph{deben} rechazarse. La \autoref{tab:rq3-modelo} y la
\autoref{fig:rq3-verdictos} resumen el desempeño por modelo.

\begin{table}[ht]
\centering
\caption{Resultados de RQ3 por modelo (150 invocaciones por modelo). Tasa de
éxito = veredictos \texttt{SUCCESS}; coincidencia con \emph{baseline} =
salida idéntica al baseline determinista; consistencia = pares caso--modelo
cuyos tres intentos coinciden. Fuente:
\texttt{output/rq3-campaign-limited/rq3-aggregated.md}.}
\label{tab:rq3-modelo}
\begin{tabular}{@{}lrrrrrr@{}}
\toprule
\textbf{Modelo} & \textbf{ÉXITO} & \textbf{INCORR.} & \textbf{ABSTEN.} & \textbf{Tasa éxito} & \textbf{Coinc.\ baseline} & \textbf{Consist.} \\
\midrule
\texttt{gpt-4o}  & 122 & 0 & 28 & 81{,}3\,\% & 78{,}7\,\% & 96{,}0\,\% \\
\texttt{gpt-4.1} & 145 & 5 & 0  & 96{,}7\,\% & 94{,}0\,\% & 98{,}0\,\% \\
\bottomrule
\end{tabular}
\end{table}

\begin{figure}[ht]
\centering
\begin{tikzpicture}
\begin{axis}[
  ybar stacked, width=0.55\linewidth, height=6cm,
  symbolic x coords={gpt-4o,gpt-4.1}, xtick=data,
  ylabel={Invocaciones (de 150)}, ymin=0, ymax=160,
  legend style={at={(0.5,-0.18)}, anchor=north, legend columns=3, font=\small},
]
\addplot+[fill=green!35] coordinates {(gpt-4o,122)(gpt-4.1,145)};
\addplot+[fill=orange!40] coordinates {(gpt-4o,28)(gpt-4.1,0)};
\addplot+[fill=red!45]   coordinates {(gpt-4o,0)(gpt-4.1,5)};
\legend{ÉXITO, ABSTENCIÓN, INCORRECTO}
\end{axis}
\end{tikzpicture}\hfill
\begin{tikzpicture}
\begin{axis}[
  ybar, width=0.40\linewidth, height=6cm,
  symbolic x coords={Elegibles,No elegibles}, xtick=data,
  ylabel={Tasa de éxito (\%)}, ymin=0, ymax=109,
  nodes near coords, every node near coord/.append style={font=\scriptsize},
  legend style={at={(0.5,-0.18)}, anchor=north, legend columns=2, font=\small},
  enlarge x limits=0.4,
]
\addplot+[fill=blue!30]  coordinates {(Elegibles,74.1)(No elegibles,100)};
\addplot+[fill=purple!30] coordinates {(Elegibles,97.2)(No elegibles,95.2)};
\legend{gpt-4o, gpt-4.1}
\end{axis}
\end{tikzpicture}
\caption{RQ3 (300 invocaciones). Izquierda: reparto de veredictos por modelo.
Derecha: tasa de éxito separando casos elegibles (deben transformarse) de no
elegibles (deben rechazarse). Datos:
\texttt{output/rq3-campaign-limited/rq3-summary.csv}.}
\label{fig:rq3-verdictos}
\end{figure}

La fotografía agregada esconde lo más interesante, que aparece al separar los
casos por elegibilidad (\autoref{tab:rq3-elegibilidad}). Los dos modelos
exhiben perfiles opuestos ante la decisión de \emph{si transformar}:

\begin{table}[ht]
\centering
\caption{Tasa de éxito de RQ3 por categoría de elegibilidad (intentos
agregados; 36 casos elegibles $\times 3$ y 14 no elegibles $\times 3$ por
modelo). Fuente: \texttt{output/rq3-campaign-limited/rq3-summary.csv}.}
\label{tab:rq3-elegibilidad}
\begin{tabular}{@{}lrr@{}}
\toprule
\textbf{Modelo} & \textbf{Elegibles ($n=108$)} & \textbf{No elegibles ($n=42$)} \\
\midrule
\texttt{gpt-4o}  & 80/108 \;(74{,}1\,\%)  & 42/42 \;(100{,}0\,\%) \\
\texttt{gpt-4.1} & 105/108 \;(97{,}2\,\%) & 40/42 \;(95{,}2\,\%) \\
\bottomrule
\end{tabular}
\end{table}

\texttt{gpt-4o} es \textbf{conservador}: no produce ni una sola transformación
incorrecta (0 \texttt{INCORRECT}) y rechaza correctamente los 42 intentos sobre
casos no elegibles (100\,\%); pero esa cautela tiene un coste, porque se
abstiene también en casos que sí debían transformarse, y su acierto sobre los
elegibles cae al 74{,}1\,\%. Sus 28 abstenciones recaen, todas, sobre casos
elegibles. \texttt{gpt-4.1} muestra el perfil \textbf{opuesto}: aplica la
transformación con acierto en el 97{,}2\,\% de los elegibles, pero a cambio se
equivoca en 5 ocasiones (las únicas \texttt{INCORRECT} de toda la campaña),
algunas de ellas transformando casos que debían rechazarse (95{,}2\,\% en no
elegibles). En términos clásicos, \texttt{gpt-4o} prioriza la precisión a costa
de la cobertura y \texttt{gpt-4.1} la cobertura a costa de algún falso
positivo. El \emph{baseline} determinista, por construcción, no incurre en
ninguno de los dos errores.

La consistencia a temperatura cero es alta en ambos modelos (96{,}0\,\% y
98{,}0\,\% de los pares caso--modelo con los tres intentos idénticos), lo que
indica que los aciertos y los fallos descritos son criterios estables de cada
modelo y no ruido de muestreo. En cuanto al coste de invocación, el tiempo de
respuesta por llamada tuvo una mediana de 1066\,ms y una media de 1649\,ms, en
un rango de 436 a 9350\,ms. Un caso resulta especialmente ilustrativo:
\texttt{ANT\_U\_R\_L\_RESOURCE\_GET\_U\_R\_L}, elegible, en el que \emph{ambos}
modelos fallaron en los tres intentos —\texttt{gpt-4o} por abstención y
\texttt{gpt-4.1} por transformación incorrecta—, lo que muestra que la
dificultad no siempre está en la elegibilidad nominal del caso. El detalle
caso por caso se conserva en
\texttt{output/rq3-campaign-limited/rq3-summary.csv}.

\subsection{Validación humana del pipeline automático}
\label{sec:validacion-humana}
El pipeline se ha construido con asistencia intensiva de herramientas de IA y,
además, sus etapas son automáticas: la detección de \lstinline|if| anidados, la
comprobación de las precondiciones P1--P5 y el cálculo de la complejidad
cognitiva. Para acotar la confianza en esos resultados se incorpora una
\textbf{validación humana} sobre una muestra aleatoria, con un principio
sencillo: tomar $N$ casos al azar, comprobar manualmente si el veredicto
automático es correcto y reportar el porcentaje de acierto por etapa.

\paragraph{Protocolo.} Se extrae una muestra de 30 casos del corpus
consolidado —20 elegibles y 10 no elegibles— mediante muestreo aleatorio
reproducible (semilla fija 42). Para cada caso, un revisor humano abre el
fichero fuente y el resultado automático y emite un veredicto
(correcto/incorrecto) en cuatro etapas: (i)~detección del patrón de
\lstinline|if| anidado; (ii)~cumplimiento o incumplimiento de las
precondiciones y su motivo de descarte; (iii)~complejidad cognitiva
\emph{before}; y (iv)~complejidad cognitiva \emph{after}. El instrumento de
recogida está en \texttt{output/rq2-human-validation/validation-sample.csv} y
el procedimiento completo en \texttt{docs/human-validation-protocol.md}.

\paragraph{Resultados.} \pendienteverif{} Pendiente de completar la inspección
manual. La \autoref{tab:validacion-humana} recoge la plantilla; el porcentaje
de acierto por etapa, con su intervalo de confianza de Wilson al 95\,\%, se
rellenará al cerrar la revisión.

\begin{table}[ht]
\centering
\caption{Validación humana del pipeline (muestra $N=30$, semilla 42).
Plantilla a completar tras la inspección manual.}
\label{tab:validacion-humana}
\begin{tabular}{@{}lccc@{}}
\toprule
\textbf{Etapa} & \textbf{Aciertos/$N$} & \textbf{\% acierto} & \textbf{IC 95\,\% (Wilson)} \\
\midrule
Detección de \lstinline|if| anidados      & --- & --- & --- \\
Precondiciones P1--P5 (y motivo)          & --- & --- & --- \\
Complejidad cognitiva \emph{before}       & --- & --- & --- \\
Complejidad cognitiva \emph{after}        & --- & --- & --- \\
\bottomrule
\end{tabular}
\end{table}

\subsection{Discusión integrada}
Leídos en conjunto, los tres resultados se refuerzan. RQ1 acota dónde la
combinación es segura y muestra que el principal obstáculo en código real es el
posible efecto colateral en la condición, no el \lstinline|else|. RQ2 confirma
que, donde es segura, la combinación reduce de forma sistemática la complejidad
cognitiva —los 104 casos elegibles mejoran—, con un efecto modesto pero
consistente y un puñado de issues de SonarQube eliminados. Y RQ3, ya sobre 50
casos, aporta el matiz más fino para responder a
\emph{«\textquestiondown{}pueden los LLM hacerlo de forma correcta y
automática?»}: los modelos saben aplicar la transformación cuando procede, pero
difieren en la frontera difícil, que no es transformar, sino \emph{abstenerse}
cuando una precondición lo prohíbe. Un modelo peca de conservador y otro de
confiado, y solo el \emph{baseline} determinista acierta en ambos extremos por
construcción. Es ahí donde el prototipo aporta su mayor valor: como referencia
fiable frente a la que medir, precisamente, esa frontera.

% =====================================================================
%  4.7 Amenazas a la validez.  ESTADO: andamiaje.
%  Sección sensible: redactar tras cerrar §4.6 (resultados).
% =====================================================================
\section{Amenazas a la validez}
\label{cap:amenazas}

\pendienteredaccion{} Redactar tras consolidar los resultados. Estructura
prevista (por tipo de validez):

\subsection{Validez de constructo.}
Proxy de complejidad cognitiva $\neq$ SonarQube oficial; validación
cruzada limitada a un subset (\autoref{sec:sonar-val}). El oráculo de RQ3
se basa en la CC y en la comparación con el baseline determinista, no en
revisión humana ni en ejecución de los tests del proyecto original.

\subsection{Validez interna.}
Determinismo del transformador validado por tests; los LLMs presentan
no-determinismo residual aun con $T=0$ (mitigado con 3 intentos). La
política \emph{one-at-a-time} limita interacciones inesperadas pero podría
ocultar efectos compuestos.

\subsection{Validez externa.}
El corpus consolidado (126 casos, 104 elegibles) procede mayoritariamente
de proyectos abiertos, pero su composición no es una muestra
estadísticamente representativa del código Java «en general»: es una
\emph{muestra intencional} de métodos con condicionales anidados (sesgo de
selección). A ello se añade una limitación de reproducibilidad: parte de
los proyectos fuente no tiene versión ni \emph{commit} anclados y el corpus
mezcla versiones de un mismo proyecto (\autoref{sec:corpus}), por lo que
reconstruirlo idéntico exigiría reanclar dichas versiones. El subset de
RQ3 ($n=50$) cubre los ejes de variación del corpus, pero sigue siendo una
muestra intencional y no una muestra aleatoria representativa del código Java
en general. Además, la comparativa de RQ3 se restringe a dos modelos del mismo
proveedor (OpenAI):
\hipotesis{} su generalización a otras familias de modelos está por
comprobar. Por último, el patrón objetivo se limita al caso base sin
\lstinline|else|/\lstinline|else if|, por lo que los resultados no son
extrapolables a refactorizaciones más complejas.

\subsection{Validez de conclusión.}
Aunque el subset de RQ3 se ha ampliado a $n=50$ casos, sigue sin sostener
inferencia estadística fuerte sobre poblaciones de código; métrica única (CC)
que puede sesgar la valoración global de calidad.

\subsection{Mitigaciones implementadas.}
Congelación del protocolo (\texttt{LlmExperimentProtocol}); trazabilidad
completa (evidencias JSON/CSV/MD por invocación); guarda de elegibilidad
v1.1 en el oráculo. \pendientemetodo{} Decidir si se amplía a otro
proveedor de LLM o se documenta la restricción de forma definitiva.

% =====================================================================
%  Capítulo 5. Conclusiones  (Bloque UMA obligatorio)
%  ESTADO: andamiaje. NO redactar todavía (depende de §4.6 y §4.7).
%  No introducir resultados nuevos aquí.
% =====================================================================
\chapter{Conclusiones}
\label{cap:conclusiones}

\pendienteredaccion{} Redactar tras cerrar resultados y amenazas.

\section{Respuesta a las preguntas de investigación}
\pendienteredaccion{} RQ1: precondiciones P1--P5 y transformación
canónica. RQ2: $\Delta$ CC agregado (re-extraído) y validación parcial con
SonarQube. RQ3: tasas de éxito y \emph{baseline match} por modelo en la
campaña fase~9. (Sin cifras hasta consolidar §\ref{sec:resultados}.)

\section{Contribuciones efectivas}
\pendienteredaccion{} Prototipo determinista, paquete de replicación,
protocolo experimental congelado y taxonomía de descartes.

\section{Limitaciones}
\pendienteredaccion{} Resumen de \autoref{cap:amenazas}.

\section{Trabajo futuro}
\pendienteredaccion{} Orientado a una posible continuación investigadora
(p.\,ej. tesis doctoral): extensión del catálogo de patrones
(\lstinline|else|/\lstinline|else if| simétricos, anidamiento triple
parcial, \emph{guards} sin efectos colaterales); validación SonarQube
completa sobre todo el corpus; ampliación y diversificación del corpus
real; comparativa con un segundo proveedor de LLM y con técnicas
\emph{few-shot}; e integración como \emph{plugin} de IDE o regla de
SonarQube.


% --- Bibliografía ---------------------------------------------------
% No se incluyen entradas inventadas. Ver bibliografia/referencias.bib.
% Estilo neutro `plain` (el formato definitivo está [PENDIENTE DE
% VERIFICACIÓN] según la guía UMA).
\bibliographystyle{splncs04}
\bibliography{bibliografia/referencias}

% --- Apéndices ------------------------------------------------------
% =====================================================================
%  Apéndices A--G.  ESTADO: andamiaje.
%  Con `\appendix`, cada \chapter se numera como «Apéndice A, B, ...».
%  El contenido fuente está en memoria/apendices/*.md (legado a migrar).
% =====================================================================
\appendix
\label{cap:apendices}

\chapter{Manual técnico}
\pendienteredaccion{} Fuente: \texttt{memoria/apendices/A-manual-tecnico.md}.

\chapter{Paquete de replicación}
\pendienteredaccion{} Fuente: \texttt{memoria/apendices/B-replicacion.md}.

\chapter{Prompt RQ3 v1.0 y oráculo}
\pendienteredaccion{} Texto literal del prompt (\texttt{LlmPromptBuilder})
y descripción del oráculo. Fuente:
\texttt{memoria/apendices/C-prompt-rq3.md}.

\chapter{Catálogo de \texttt{DiscardReason}}
\pendienteredaccion{} Diez categorías (verificadas en
\texttt{DiscardReason.java}): \texttt{OUTER\_HAS\_ELSE},
\texttt{OUTER\_BLOCK\_MULTIPLE\_STATEMENTS},
\texttt{SINGLE\_STATEMENT\_NOT\_IF}, \texttt{INNER\_HAS\_ELSE},
\texttt{METHOD\_CALL\_IN\_CONDITION}, \texttt{ASSIGNMENT\_IN\_CONDITION},
\texttt{INCREMENT\_OR\_DECREMENT\_IN\_CONDITION},
\texttt{LAMBDA\_IN\_CONDITION}, \texttt{OBJECT\_CREATION\_IN\_CONDITION},
\texttt{NO\_NESTED\_IF\_PATTERN}. Fuente:
\texttt{memoria/apendices/D-discard-reasons.md}.

\chapter{Tablas extendidas RQ2 / RQ3}
\pendienteredaccion{} Fuente: \texttt{memoria/apendices/E-tablas-extendidas.md}.

\chapter{Procedimiento SonarQube}
\pendienteredaccion{} Fuente: \texttt{memoria/apendices/F-sonarqube.md} y
\texttt{docs/}.

\chapter{Trazabilidad RQ $\leftrightarrow$ artefactos}
\pendienteredaccion{} Fuente: \texttt{memoria/apendices/G-trazabilidad.md}.


\end{document}
